%Sample for icsthesis.cls
%Author: ATMatias
%Reminders: Omit \lipsum commands

\documentclass{icsthesis}
% APA-style citations and reference list: icsthesis loads apacite [natbibapa].
% Do not use the cite package here—it forces numeric [n] citations and conflicts with APA author–year format.
\usepackage[T1]{fontenc}
\usepackage{graphicx,amssymb,fancyhdr,setspace,xurl,xltabular}
\usepackage{microtype}
\setlength{\emergencystretch}{3em}
% Help long \path{...} strings break at delimiters (scikit-learn calls, file paths).
\makeatletter
\g@addto@macro\UrlBreaks{\do\_\do\.\do\/\do\-\do\:\do\=\do\+\do\;\do\'}
\makeatother

%\usepackage{showframe} %for debugging the margins
                              
                              
%Replace the values here, various commands depend on this variables                                                                                                                                                            
\renewcommand{\TITLE}{EXPLAINABLE ARTIFICIAL INTELLIGENCE FOR FINANCIAL PRODUCT RECOMMENDATIONS IN PHILIPPINE FINANCIAL TECHNOLOGY PLATFORMS}
\renewcommand{\AUTHOR}{ASHTON STEPHONIE MATIAS}
\renewcommand{\DEGREE}{BACHELOR OF SCIENCE}
\renewcommand{\MAJOR}{Computer Science}
\renewcommand{\MONTH}{MAY}
\renewcommand{\YEAR}{2026}

\begin{document}
	
	\begin{frontmatter}
		%create title
		\maketitle
				
		\begin{approvalpage}
			The special problem hereto attached entitled \TITLE , prepared and submitted by \AUTHOR , in partial fulfillment of the requirements for the degree of \DEGREE\ (\MAJOR) is hereby accepted.
			
			\begin{multicols}{2}
				\centering
				\addsignature{FERMIN ROBERTO G. LAPITAN}{Member, Guidance Committee} %custom command for adding signature field
				\columnbreak
				\addsignature{JOHN PATRICK VJ. ALBACEA}{Member, Guidance Committee}
			\end{multicols}
            \begin{multicols}{2}
				\centering
				\addsignature{JUAN MIGUEL J. BAWAGAN}{Member, Guidance Committee} %custom command for adding signature field
				\columnbreak
				\addsignature{RACHEL EDITA O. ROXAS}{Adviser}
			\end{multicols}
			
			Accepted in partial fulfillment of the requirements for the degree of \DEGREE\ (\MAJOR). 
			\addsignature{RACHEL EDITA O. ROXAS}{Director, Institute of Computer Science}
			\addsignature{CHRYSLINE MARGUS N. PI\~{N}OL}{Dean, Graduate School\\University of the Philippines Los Ba\~{n}os}
		\end{approvalpage}
		
		\begin{biosketch}
			Ashton Stephonie T. Matias is a Bachelor of Science in Computer Science student at the University of the Philippines Los Ba\~{n}os. He has interests in artificial intelligence, machine learning, data analytics and data science for business intelligence. He is driven by a commitment to developing inclusive and accessible solutions for marginalized communities in the Philippines.

            Beyond academics, Ashton actively participates in student leadership and extracurricular activities, particularly in the Alliance of Computer Science Students (ACSS) UPLB and university athletics. He is also a varsity student-athlete for the UPLB Men's Futsal Team and a scholar dedicated to balancing academic excellence, leadership, and professional growth.
			
			\addauthorsignaturefield
		\end{biosketch}	
		
		%acknowledgement
		\begin{acknowledgement}
			This work would not exist without the people who made the author believe they could.

            To God, who set the stars in place and this path before him -- all glory belongs to You.

            To his adviser, Dr. Rachel Edita O. Roxas, whose guidance was the compass in every uncertain chapter -- thank you for seeing the finish line even when it was not yet visible.

            To his family, the roots that held firm through every storm, your love was never loud, but it was always there.

            To his friends, \textit{Dove Girlies}, who turned the hardest seasons into something worth remembering -- thank you for the warmth and laughter in times of despair.

            To \textit{v}, who stayed through the late nights and the self-doubt, who believed in him on the days he forgot how to -- this one is for you.


            And to everyone who ever said \textit{you can do this} -- he did. And he carried every one of you with him to the end.
		\end{acknowledgement}
		
		%TOC
		\maketableofcontents
		
		%LOT
		\makelistoftables

		%LOF
		\makelistoffigures
	
		%ABSTRACT
		
		\begin{abstractwithpageno}	
		\\
		\textbf{\AUTHOR}, University of the Philippines Los Ba\~{n}os, \MONTH\ \YEAR. \textbf{\TITLE}
		\\Major Professor: Dr. Rachel Edita O. Roxas\\
			
			Financial inclusion in the Philippines remains limited even as digital platforms such as GCash, Maya, and BPI expand, yet many fintech recommenders offer no intelligible rationale for their suggestions. This study builds FinXplain: a synthetic dataset of 5{,}000 Filipino profiles (16 features) grounded in PSA FIES 2023, PSA LFS 2023, the PSA 2020 Census, and BSP FIS 2021, with rule-based labels for eight savings and investment products from those providers. An Explainable Boosting Machine (EBM) is the primary model; XGBoost is a benchmark. Both are evaluated on a stratified 20\% holdout using accuracy and weighted precision, recall, F1, and one-vs-rest AUC-ROC. A Next.js and FastAPI prototype exposes the EBM's native shape functions, plain-language reasons, and counterfactuals; SHAP and LIME are used only offline against XGBoost. The EBM achieves weighted F1 and AUC-ROC of $0.8711$ and $0.9761$, beating XGBoost by $5.93$ and $0.90$ percentage points and clearing a two-point parity bar on both, with five-fold CV accuracy $0.870 \pm 0.006$. The results support glass-box models for transparent, regulator- aligned savings guidance without giving up predictive performance.

			\vspace{2ex}
			\noindent\textbf{Keywords:} counterfactual explanations, explainable artificial intelligence, explainable boosting machine, LIME, Philippine financial technology, savings product recommendation, SHAP, synthetic data generation, XGBoost
		\end{abstractwithpageno}

	\end{frontmatter}
	
	%BODY OF THESIS HERE (Use up to 3 levels only, (sec -> subsec -> subsubsec)
	\begin{mainmatter}
		\section{INTRODUCTION}
			\subsection{Background of the Study}
				Financial technology (fintech) has rapidly transformed the financial services landscape in the Philippines. Mobile financial technology adoption continues to surge. According to data from a report \cite{ref1}, nearly 80\% of Filipinos aged 15 and above are expected to use mobile financial technology services by the end of 2024. This growth is not limited to payments; digital banking, lending applications, and personal finance platforms are also expanding as digital financial ecosystems mature \cite{ref2}. The Bangko Sentral ng Pilipinas (BSP), recognizing this momentum, has adopted forward-looking regulatory strategies such as the ``test-and-learn'' approach to support financial technology innovation and broaden financial inclusion \cite{ref3}.

                Despite these advancements, significant challenges persist, particularly in financial literacy. Studies from the Philippine Institute for Development Studies (PIDS) note that low digital literacy and limited technical know-how hinder consumer adoption of more sophisticated financial services \cite{ref4}. Financial literacy among Filipinos remains low: only 25\% of Filipino adults were found to be financially literate, placing the Philippines 13th out of 21 developing Asian economies \cite{ref5}. These gaps in understanding make users more dependent on basic financial tools and less confident in interpreting or evaluating complex financial recommendations.

                As financial technology applications expand their reach, many have begun incorporating artificial intelligence (AI) and machine learning to deliver personalized financial insights. For example, GCash employs GCoach AI to provide tailored suggestions on saving, investing, and money management \cite{ref6}. Similarly, BPI integrates Track \& Plan, an AI-enabled feature that analyzes transaction histories to identify spending patterns and suggest opportunities such as investments when excess cash is detected \cite{ref7}. These capabilities reflect a broader trend in financial technology innovation: enhancing user engagement through personalized and automated financial guidance.

                However, these AI-driven recommendation tools function as ``black boxes''; users are presented with recommendations without understanding the reasoning that produced them. The lack of explainability limits user trust and may discourage adoption, especially in financial contexts where understanding and confidence are essential for responsible decision-making. Furthermore, the development of transparent and contextually relevant AI systems in the Philippines is constrained by the limited availability of public financial datasets. Since real financial technology transaction data are protected by privacy laws and industry confidentiality, researchers must often rely on synthetic datasets that simulate Filipino financial behaviors while safeguarding privacy.
                
			\subsection{Statement of the Problem}
				Financial technology platforms in the Philippines are increasingly using artificial intelligence to provide users with personalized financial product recommendations. However, these systems typically operate in a manner that lacks transparency, which means that users do not receive clear explanations for why specific products are suggested to them. This lack of explainability can weaken user trust, limit financial understanding, and reduce the willingness of consumers to adopt or act on recommended financial products. Although these AI-driven recommendation systems use behavioral and financial data to generate predictions, the reasoning behind their outputs often remains inaccessible to users. 

                In addition, the development of locally grounded recommendation systems is challenged by the limited availability of publicly accessible Philippine financial datasets. Because financial technology user data are protected by privacy, security, and regulatory constraints, researchers cannot easily obtain real-world transactional data. As a result, there is a need to construct synthetic datasets that accurately simulate Filipino financial behaviors while enabling safe experimentation and system prototyping.

                To address this gap, the present study examines how explainable AI techniques can be incorporated into a Philippine fintech savings-product recommender, and how transparent explanations can support a more responsible and user-centered design that reinforces financial literacy in the Philippine financial-technology landscape.
                
            \subsection{Objectives of the Study}
				This study aims to design and develop a web-based prototype financial-product recommender system that provides transparent, plain-language explanations for its recommendations Specifically, it aims to:

                \begin{enumerate}
                    \item Generate a synthetic dataset of Filipino fintech user profiles across 16 features, with distributions parameterized from published Philippine statistical sources (PSA FIES, PSA LFS, PSA 2020 Census, and BSP FIS).
                    \item Label the dataset through rule-based conditional logic aligned to the eligibility criteria of the eight selected savings and investment products from GCash, Maya, and BPI.
                    \item Train an Explainable Boosting Machine (EBM) and an XGBoost benchmark, evaluated on a stratified 20\% held-out test set using accuracy, weighted precision, weighted recall, weighted F1-score, and weighted one-vs-rest (OvR) AUC-ROC.
                    \item Show that EBM attains a weighted F1-score and AUC-ROC within two percentage points of, or greater than, XGBoost, demonstrating that native interpretability can be retained without a substantive loss in predictive performance. The two-percentage-point threshold follows the performance-parity convention used in prior interpretable-vs.-black-box benchmarking on tabular data, where differences below this magnitude are generally treated as substantively equivalent \cite{ref15}, \cite{ref50}, \cite{ref51}. 
                \end{enumerate}
                
            \subsection{Significance of the Study}
				Financial recommendations increasingly shape how Filipinos save, spend, borrow, and invest. In finance, where decisions affect personal well-being, explainability becomes just as important as prediction quality.

                For the financial technology industry, the prototype provides a concrete reference architecture for building more transparent AI-driven product-recommendation features that align with the expectations of Filipino consumers and with BSP and Data Privacy Act requirements for intelligible automated decisioning. Rather than claiming measurable gains in user experience or in product adoption---neither of which has been evaluated here---the study is positioned as demonstrating a technically feasible, regulation-aligned pattern that providers such as GCash, Maya, and BPI, or academic groups working on Philippine fintech, can adapt and empirically validate in subsequent work.

                On a societal level, the study aims to contribute to user empowerment in financial decision-making. When users are shown why a system recommends a product, the intent is that they gain insights into their financial behaviors and constraints, supporting more informed choices and more inclusive and responsible financial ecosystems in the Philippines; measurable effects on behavior and adoption are left to follow-up empirical work.

                For the academic community, the study contributes a reusable Philippine-contextualized synthetic fintech dataset parameterized jointly from PSA FIES \cite{ref40}, PSA LFS \cite{ref41}, the PSA 2020 Census \cite{ref39}, and BSP Financial Inclusion Survey \cite{ref35} sources. To the best of the author's knowledge no prior publicly available synthetic dataset jointly parameterizes these four sources for the Philippine retail savings market, and the dataset, generation script (fixed seed~\cite{ref60}), and feature-to-source mapping are intended as a research artifact for subsequent Philippine fintech ML work.
                
            \subsection{Scope and Delimitations}
				\subsubsection{Scope}
                    This study develops a research prototype, not a production-grade financial-technology application. The recommendation engine covers eight savings and investment products from three Philippine fintech providers: GCash (GSave, GInvest), Maya Bank (Savings, Personal Goals, Time Deposit Plus), and BPI (Regular Savings, \#SaveUp, Plan Ahead Time Deposit). The training corpus is a synthetic dataset of 5{,}000 user profiles whose distributions are parameterized from published Philippine statistical sources (PSA FIES, PSA LFS, PSA Census, BSP FIS). The recommendation engine uses an Explainable Boosting Machine (EBM) as the primary model and XGBoost as a benchmark comparator; explanations are produced by the EBM's native additive shape functions, with SHAP, LIME, and counterfactual methods used for offline benchmarking. A web-based interface collects 16 user inputs through a three-step wizard and returns a predicted product together with top-3 alternatives, a top-5 key-factors bar chart, plain-language supporting and opposing reasons, and counterfactual suggestions.

                \subsubsection{Delimitations}
                    The study does not address regulatory compliance in depth beyond citing the enabling statutes (RA 10173, RA 11765) and their operational implications. It does not handle structured credit products (personal loans, mortgages, credit cards) or complex investment products (stocks, bonds, mutual funds outside UITFs). It does not use real financial records, nor does it integrate with any live provider API. The current input schema captures stated suitability preferences (savings goal, risk tolerance, investment horizon) and a snapshot financial profile, but does not capture \emph{behavioural} or \emph{ongoing-suitability} signals---recurring-deposit patterns, withdrawal frequency, balance volatility, and similar longitudinal transactional behaviours that a real Philippine advisory engagement would normally incorporate. It does not conduct a user study, so usability claims rest on design principles rather than empirical validation. Long-term behavioral outcomes (e.g., whether a displayed explanation actually changes financial behavior) are not measured.
		
		\section{REVIEW OF LITERATURE}
		Recommendation systems have become central to how users navigate digital services, with applications extending from e-commerce and entertainment into financial technology \cite{ref13}. Financial recommenders differ from entertainment recommenders in three material respects: they operate over sensitive personal data, they are subject to financial-consumer-protection regulation, and the cost of a poor recommendation is financial loss rather than inconvenience \cite{ref13}, \cite{ref14}. These differences motivate the present study's emphasis on explainability and on a Philippines-specific regulatory and behavioral framing.

			\subsection{Financial Technology}
				Financial technology has transformed traditional financial services by integrating digital systems, automation, and data-driven processes that enable more efficient payments, credit access, and financial management. Globally, financial technology adoption continues to accelerate, driven by advancements in mobile banking, electronic payments, digital lending, and alternative finance platforms \cite{ref9}. Benefits include increased accessibility, reduced transaction costs, and personalized financial tools, although challenges such as cybersecurity, digital literacy, and uneven regulatory capacity remain.

                In the Philippines, financial technology adoption has significantly increased, especially during and after the COVID-19 pandemic. Quimba, Barral, and Carlos \cite{ref9} noted strong growth in digital payments and mobile wallets, but also identified structural constraints such as low financial literacy, inadequate talent supply, and inconsistent regulatory frameworks. These findings highlight the need for transparent and trustworthy digital financial systems that users not only adopt but also understand. At the national level, the principal publicly available Philippine household financial datasets are the Philippine Statistics Authority Family Income and Expenditure Survey (FIES) \cite{ref40}, the PSA Labor Force Survey (LFS) \cite{ref41}, and the Bangko Sentral ng Pilipinas Financial Inclusion Survey (FIS) \cite{ref35}; these are the primary sources used to parameterize the synthetic population in the Methodology chapter.

			\subsection{Data Science in Financial Technology}
				Data science is central to modern financial technology innovation because financial decisions increasingly rely on large volumes of transactional and behavioral data. According to Provost and Fawcett \cite{ref10}, high-stakes analytics systems require rigorous data pipelines consisting of dataset construction, cleaning, and preprocessing to ensure reliability and fairness. In financial environments, data quality is particularly critical because noise, missing values, or outliers can distort predictions and risk assessments.

                Financial datasets often undergo extensive data cleaning, such as removing duplicates, imputing missing values, and normalizing skewed distributions. These processes reduce bias and improve model stability \cite{ref10}. Data preprocessing techniques, including encoding categorical variables and scaling numerical features, are essential to enable machine learning algorithms to correctly interpret user financial behaviors \cite{ref11}.

			\subsection{Recommendation Systems in Financial Technology}
				Recommendation systems are increasingly used to personalize financial services such as loan offers, investment bundles, savings plans, and credit products. These systems rely on user profiles, transaction histories, and behavioral data to tailor financial recommendations, improving both relevance and engagement \cite{ref13}. Personalized recommendations help users navigate complex financial choices, but they also raise challenges including data privacy, fairness, accuracy, and lack of transparency.

                In financial technology, recommendation engines must operate under high-stakes conditions, where poor suggestions could negatively affect savings, creditworthiness, or investment performance. Thus, ensuring accountability and explainability is as important as optimizing predictive performance. As Ricci et al. explain \cite{ref13}, opaque recommendation algorithms can undermine user trust, especially when decisions affect financial well-being. This reinforces the need for explainable and interpretable recommendation systems in financial applications.

			\subsection{Machine Learning in Financial Technology}
				Machine learning models play a core role in powering financial technology recommendation systems. Supervised learning algorithms, such as logistic regression, decision trees, gradient boosting, and neural networks, are widely used for credit scoring, customer segmentation, and product recommendation \cite{ref11}, \cite{ref14}, \cite{ref56}. Their ability to capture nonlinear patterns and complex interactions makes them highly effective in predicting financial behaviors and user preferences.

                However, ML systems deployed in financial settings must meet requirements that go beyond accuracy. They must be scalable, robust to noisy data, and compliant with fairness and transparency standards \cite{ref14}. As models grow more complex, explaining their decisions becomes increasingly challenging, which has intensified interest in interpretable and explainable AI techniques in the financial domain.

			\subsection{Explainable Boosting Machine and XGBoost}
				The Explainable Boosting Machine (EBM) is an interpretable machine-learning model built on generalized additive models with pairwise interactions (GA$^2$M). The underlying model class and training algorithm were introduced by Lou et al.\ \cite{ref48}, \cite{ref49}, and the open-source InterpretML implementation used in this study is the one described by Nori et al.\ \cite{ref15}. EBM learns a shape function for each feature, allowing users to visualize how factors such as income, savings behavior, or spending patterns contribute to a recommendation. Its transparency makes EBM suitable for applications that require human understanding of model decisions, including healthcare, finance, and risk management.

                By contrast, XGBoost is a high-performance gradient boosting algorithm developed by Chen and Guestrin \cite{ref16}. It remains a state-of-the-art model for many Kaggle competitions and financial prediction tasks because it supports parallel computation, handles missing values efficiently, and provides exceptional predictive accuracy. Although it is less interpretable than EBM, post-hoc explainability tools can recover meaningful explanations.

                Together, EBM and XGBoost enable a balanced approach in financial technology recommenders. EBM prioritizes interpretability, while XGBoost provides a performance benchmark for comparison.

			\subsection{Explainable Artificial Intelligence Methods}
				Explainability is crucial in financial AI systems due to regulatory requirements, ethical considerations, and user trust. Financial decisions can significantly affect individuals, so regulatory bodies increasingly demand transparency in automated systems, especially in credit, lending, and personalized recommendations \cite{ref17}. Explainability helps identify biases, justify model predictions, enhance accountability, and provide users with understandable reasons for recommendations.

				\subsubsection{Shapley Additive Explanations}
					Based on cooperative game theory, SHAP (SHapley Additive exPlanations) assigns each feature a Shapley value representing its contribution to a prediction \cite{ref18}. SHAP is widely used in finance because it provides both global feature importance and local explanations for individual users.

				\subsubsection{Local Interpretable Model-Agnostic Explanations}
					Introduced by Ribeiro et al. \cite{ref19}, LIME (Local Interpretable Model-agnostic Explanations) approximates a complex model locally with an interpretable surrogate model. This helps users understand which features influenced a recommendation for their specific case.

				\subsubsection{Counterfactual Explanations}
					Counterfactuals indicate how minimal changes in input features can alter the model output, such as ``If you saved PHP 1,000 more per month, you would qualify for a higher-tier product.'' Wachter, Mittelstadt, and Russell \cite{ref20} note that counterfactuals are particularly valuable because they provide actionable, human-centered guidance.

			\subsection{Synthetic Data in Machine Learning Research}
				Synthetic data refers to artificially generated records that mimic the statistical properties of real-world data without containing information about actual individuals. Jordon et al. \cite{ref83} define synthetic data in the machine learning context as data generated by an algorithmic process that preserves the distributional and relational structure of an original (or notional) population while decoupling the outputs from any identifiable record. The motivations for using synthetic data fall into four commonly cited categories: (i) privacy preservation, where the release of real records would violate data-protection law or ethical norms; (ii) data scarcity, where real datasets are too small, expensive, or restricted for reliable model development; (iii) class imbalance correction, where minority classes are under-represented in real data; and (iv) controlled experimentation, where researchers must fix the data-generating process to isolate model behavior from sampling noise \cite{ref83}, \cite{ref46}.

                Synthetic data is particularly relevant for Philippine fintech research because publicly available Philippine household financial datasets are scarce. The Philippine Data Privacy Act of 2012 (Republic Act No. 10173) \cite{ref45} restricts the use of real financial records for academic research without explicit consent and institutional data-sharing agreements, and BSP consumer data is not made available for unrestricted academic analysis. In the absence of open Philippine fintech micro-data, synthetic data grounded in published statistical sources is a defensible methodological choice \cite{ref46}, \cite{ref83}.

			\subsection{Recommendation Systems in Personal Finance}
				The recommendation-systems literature in fintech has focused predominantly on credit scoring and loan underwriting, where a binary ``approve/reject'' label is directly observable from lender records \cite{ref56}, \cite{ref57}. Personal-finance product recommendation (choosing among savings accounts, time deposits, and investment products based on a user's goals, risk tolerance, and horizon) is comparatively under-studied because product-suitability labels are not routinely captured in transaction logs. Ricci et al.\ \cite{ref13} note that in high-stakes personal-finance settings, the accuracy-trust trade-off is asymmetric: a highly accurate but opaque recommendation is less valuable than a slightly less accurate but transparent one, because the user's ability to act on (and contest) the recommendation depends on understanding why it was made.

                For quantitative context, West's \cite{ref56} neural-network credit-scoring work on the UCI German Credit (1{,}000 applicants) and Australian Credit (690 applicants) benchmark datasets reports classification accuracies typically in the 0.70--0.85 range on \emph{binary} ``good vs.\ bad credit'' labels with real outcomes, depending on the network architecture and resampling protocol; the original paper reports these figures as accuracies rather than weighted F1 scores, and the exact values are conditional on the specific dataset split and network configuration reported in West's Tables~2--5. Bussmann et al.\ \cite{ref14} report 0.60--0.75 accuracy on real European credit-risk data using explainable tree ensembles. Translating these benchmarks to the present study is imperfect because the task here is 8-class rather than binary and the labels are synthetic, but they establish the performance band that a reasonable interpretable tabular classifier can be expected to occupy on this class of problem. No prior study reports multi-class Philippine fintech product-recommendation metrics on a comparable task; filling this comparability gap is one contribution of the present study.

			\subsection{Behavioral Finance Foundations for Product-Suitability Assessment}
				The feature space of a product-suitability recommender is not a neutral choice. The features that matter (age, income, dependents, risk tolerance, horizon, and goal) are motivated by the behavioral-finance literature rather than by pure predictive power. Modigliani and Brumberg's life-cycle hypothesis \cite{ref65} predicts that an individual's savings and consumption pattern co-varies systematically with age and expected lifetime income, which motivates the treatment of \emph{age} and \emph{num\_dependents} as first-class features in the generator. Kahneman and Tversky's prospect theory \cite{ref62} motivates the discrete \emph{risk\_tolerance} variable rather than a continuous risk score, because prospect-theoretic decision-making is more naturally captured by ordinal risk attitudes than by a point estimate of utility curvature. Lusardi and Mitchell's \cite{ref43} work on financial literacy and planning connects the \emph{investment\_horizon} and \emph{savings\_goal} features to observed differences in retirement readiness and savings-product choice, and Browning and Lusardi's \cite{ref55} household-saving framework connects the \emph{disposable\_income}, \emph{expense\_ratio}, and \emph{savings\_to\_income\_ratio} derived features to theoretical saving propensities. These behavioral-finance foundations are imported into the synthetic-data generator through the conditional-sampling rules and are revisited in the circular-reasoning discussion because they are what makes a rule-based label defensible as a proxy for real product-suitability outcomes.

			\subsection{Synthesis and Research Gap}
				Three gaps emerge from the literature surveyed above and together motivate the present study. First, Philippine fintech product recommendation has not been studied as a supervised machine-learning task; the existing Philippine fintech literature is dominated by descriptive surveys \cite{ref9}, \cite{ref35} rather than modeling studies. Second, the Explainable Artificial Intelligence (XAI) literature has concentrated on post-hoc methods (SHAP, LIME) applied to black-box models \cite{ref17}, \cite{ref18}, \cite{ref19}, but the growing intelligible-models literature \cite{ref15}, \cite{ref48}, \cite{ref49}, \cite{ref50}, \cite{ref51} argues that on structured tabular data the more defensible path is to use an inherently interpretable model and to evaluate it against a strong black-box benchmark---a comparison that has not been reported for Philippine fintech product recommendation. Third, the synthetic-data-in-finance literature \cite{ref46}, \cite{ref83} defends parametric generation under privacy constraints but does not provide a Philippines-grounded blueprint that combines PSA and BSP sources with behavioral-finance-motivated feature choices. The present study addresses these three gaps jointly: a Philippines-grounded synthetic dataset, an EBM-vs-XGBoost head-to-head on a licensed 8-product catalogue, and a production path in which user-facing explanations are read directly from the model rather than from a post-hoc approximation.

		
		\section{METHODOLOGY}
		This section presents the methodology workflow used for this study as shown in Fig.1. The host machine, development tools, and dataset used for this study are discussed in detail. Initial stages include data cleaning and data preprocessing in preparation for training the recommendation model. After training, the XAI methodologies and prototype features are outlined.

			\begin{figure}[ht]
			\begin{center}
				\vspace{4ex}
				\includegraphics[width=0.5\textwidth]{workflow.png}
				\label{fig:workflow}
				\caption{Methodology workflow for the study.}
			\end{center}
			\end{figure}

			\subsection{Host Machine}
				\begin{itemize}
				    \item \textbf{Operating System:} Windows 11 Home 64-bit
				    \item \textbf{Processor:} AMD Ryzen 5 8645HS w/ Radeon 760M Graphics (up to 5.00 GHz, 6 cores, 12 threads)
				    \item \textbf{Memory:} 16GB DDR5
				\end{itemize}

			\subsection{Development Tools}
				\subsubsection{Backend Technologies}
					FastAPI is a modern, high-performance Python web framework designed for building APIs. It supports asynchronous processing, automatic documentation generation (Swagger UI), and strong input validation with Pydantic \cite{ref21}. Its efficiency makes it suitable for ML-based applications requiring fast response times.

				\subsubsection{Frontend Technologies}
					The prototype frontend is implemented with Next.js~14 (using the App Router) and React, which together provide server-rendered and client-interactive pages for the user-facing interface \cite{ref22}. Styling uses Tailwind CSS, animations use Framer Motion, and state persistence across pages uses Zustand.

				\subsubsection{Model Libraries}
					The Explainable Boosting Machine is trained with InterpretML's \emph{ExplainableBoostingClassifier} \cite{ref15}. The XGBoost benchmark is trained with the \emph{XGBClassifier} API \cite{ref16}. All preprocessing, train/test splitting, and evaluation metrics use scikit-learn \cite{ref74}. SHAP \cite{ref18}, LIME \cite{ref19}, and counterfactual analyses \cite{ref20} are implemented in Python for offline model analysis.

			\subsection{Financial Products Included in the Prototype}
				The FinXplain prototype recommends savings and investment products from three Philippine fintech providers: GCash, Maya (Maya Bank), and BPI (Bank of the Philippine Islands). These three providers were selected because they collectively represent the three dominant categories of Philippine financial service providers accessible to retail consumers: mass-market electronic-wallet platforms, BSP-licensed digital banks, and universal banks. This three-provider framework captures the structural diversity of the Philippine financial services market reported by the Bangko Sentral ng Pilipinas \cite{ref35} and by the Philippine fintech landscape analysis of Quimba, Barral, and Carlos \cite{ref9}.

				\subsubsection{GCash}
					GCash is the dominant electronic-wallet platform in the Philippines, holding approximately 89\% of the national mobile-wallet market share and serving over 94 million registered users as of 2025 \cite{ref9}, \cite{ref35}. Its zero-cost account creation, smartphone-first delivery, and nationwide cash-in/cash-out network make GCash the primary digital financial access point for unbanked and underserved Filipinos \cite{ref35}, \cite{ref9}. GCash is regulated by the BSP as an electronic-money issuer and is supervised under the same consumer-protection framework that applies to banks \cite{ref35}. Two GCash products are included in the prototype: \emph{GSave}, a savings account powered by CIMB Bank Philippines and accessible through the GCash application, and \emph{GInvest}, an investment marketplace that gives users access to Unit Investment Trust Funds (UITFs). GSave is selected because it has no minimum balance requirement and is Philippine Deposit Insurance Corporation (PDIC)-insured, making it an appropriate low-barrier entry point for first-time savers \cite{ref35}, \cite{ref66}. GInvest is selected because its PHP 50 minimum investment lowers the traditional entry threshold for UITFs and opens market-linked investment to risk-tolerant users with long horizons \cite{ref9}, \cite{ref66}. Together, these two products cover the low-barrier savings and accessible investment segments, respectively.

				\subsubsection{Maya (Maya Bank)}
					Maya Bank is one of only six BSP-licensed digital banks in the Philippines, chartered under BSP Circular No.~1105, Series of 2020 \cite{ref36}. Its savings rates are well above the market average of traditional universal banks \cite{ref35}, \cite{ref9}, which makes Maya financially relevant for the moderate-income digital-savvy segment that the study population includes. Three Maya products are represented: \emph{Maya Savings}, a high-interest liquid savings account with no minimum balance; \emph{Maya Personal Goals}, a goal-based savings product that allows users to create up to five named goal accounts with tiered interest rates; and \emph{Maya Time Deposit Plus}, a digital time-deposit product with 3-, 6-, and 12-month tenors. These three products jointly cover the liquid-digital-savings segment, the goal-based savings segment that draws on the behavioral-finance literature on commitment mechanisms \cite{ref63}, and the locked-term deposit segment for users willing to trade liquidity for higher guaranteed returns. Maya deposits are insured by the Philippine Deposit Insurance Corporation up to PHP 1,000,000 per depositor, providing the same statutory protection as traditional bank deposits \cite{ref36}.

				\subsubsection{BPI (Bank of the Philippine Islands)}
					BPI is the second-largest universal and commercial bank in the Philippines by total assets \cite{ref37}. It is also the oldest bank in Southeast Asia, established in 1851, and holds a BSP universal banking license with PDIC membership \cite{ref37}. BPI represents the traditional universal-banking segment, which remains essential for formally employed and risk-averse Filipino users who require payroll accounts, physical branch access, and the institutional trust associated with a long-established universal bank \cite{ref35}, \cite{ref37}. Three BPI products are included: \emph{BPI Regular Savings}, a traditional passbook/debit savings account; \emph{BPI~\#SaveUp}, a digital auto-save account operated through the BPI mobile application; and \emph{BPI Plan Ahead Time Deposit}, a long-term five-year fixed-rate time deposit. These three products cover conservative liquid savings, digital goal-based saving, and long-horizon deposits, respectively. Including a traditional universal bank alongside GCash and Maya was necessary to complete the product-coverage spectrum, because roughly 23\% of Filipino adults still rely primarily on traditional bank accounts even after the rapid expansion of electronic-money accounts \cite{ref35}, \cite{ref66}.

				\subsubsection{Product Coverage}

					Table~\ref{tab:products} summarizes the eight products included in the prototype, their providers, product type, posted interest rate, minimum balance, liquidity, and risk level.

					\begin{table}[ht]
					\vspace{4ex}
					\centering
					\caption{Summary of Financial Products Included in the FinXplain Prototype}
					\label{tab:products}
					\small
					\begin{tabular}{|p{1.3cm}|l|p{1.7cm}|p{1.5cm}|l|p{1.4cm}|p{1.2cm}|}
						\hline
						\textbf{Product} & \textbf{Provider} & \textbf{Type} & \textbf{Interest Rate} & \textbf{Min. Balance} & \textbf{Liquidity} & \textbf{Risk}\\
						\hline
						GSave (CIMB)             & GCash & Digital Savings        & 2.6\% p.a.                 & PHP 0      & High & Low         \\
						\hline
						GInvest (UITFs)          & GCash & Investment             & Variable (market)          & PHP 50     & Medium & Medium\newline--High \\
						\hline
						Maya Savings             & Maya  & Digital Savings        & 3.0\% base, up to 15\%     & PHP 0      & High & Low         \\
						\hline
						Maya Personal Goals      & Maya  & Goal-Based Savings     & 4--8\% tiered              & PHP 0      & High & Low        \\
						\hline
						Maya Time Deposit Plus   & Maya  & Digital Time Deposit   & 3.5--6.0\% p.a.            & PHP 5,000  & Low (locked) & Low   \\
						\hline
						BPI Regular Savings      & BPI   & Traditional Savings    & 0.0625\% p.a.              & PHP 3,000  & High & Low         \\
						\hline
						BPI\newline~\#SaveUp             & BPI   & Digital Auto-Save      & 0.0925\% p.a.              & PHP 1      & High & Low         \\
						\hline
						BPI Plan Ahead TD        & BPI   & Time Deposit           & Fixed (branch inquiry)     & PHP 50,000 & Low (5-yr lock) & Low\\
						\hline
					\end{tabular}
					\\[2pt]
					{\footnotesize Note. Interest rates verified April 2026 and subject to change. Sources: official GCash, Maya Bank, and BPI product documentation, synthesized in \cite{ref35}.}
					\vspace{4ex}
					\end{table}

				\subsubsection{Academic Precedent for Product Selection Criteria}
					Financial-product suitability assessment in this study draws on three strands of literature. First, work on financial literacy and household financial decision-making shows that product choice should be matched to income, dependents, risk tolerance, and time horizon rather than selected purely on headline interest rates \cite{ref43}, \cite{ref55}. Lusardi and Mitchell \cite{ref43} document that mismatches between product complexity and financial literacy degrade welfare, motivating the inclusion of both low-barrier and higher-commitment products in the recommendation space. Second, the recommender-systems literature emphasizes that product diversity and catalog coverage are essential for fair and useful recommendations; a system that recommends only a narrow band of products systematically under-serves segments of the user population \cite{ref13}. The eight-product catalog used here spans liquid digital savings, goal-based savings, digital and traditional time deposits, traditional savings accounts, and market-linked investment, covering the full spectrum of retail savings/investment needs recognized in Philippine personal finance. Third, the requirement for explainability in this domain is grounded in Philippine law: the Data Privacy Act (Republic Act No.~10173) \cite{ref45} mandates that data subjects be informed of the logic involved in automated decision-making, and the Financial Consumer Protection Act (Republic Act No.~11765) \cite{ref44} obliges financial service providers to disclose material information in a form that consumers can reasonably understand. These statutes constitute the regulatory basis for the study's explainability requirement and, in turn, for its choice of an inherently interpretable primary model.

			\subsection{Dataset Source and Synthetic Data Generation}
				Because transactional datasets from Philippine financial technology providers are not publicly available due to data-privacy regulations, confidentiality agreements, and bank-level security protocols \cite{ref45}, \cite{ref46}, the study generates a synthetic dataset that approximates the financial behavior of Filipino digital-finance users. A fixed pseudorandom seed (\emph{np.random.seed(42)}) is set so that the entire dataset is deterministically reproducible \cite{ref60}.

				\subsubsection{Sample Size}
					The synthetic dataset contains $n = 5{,}000$ observations, each described by 16 input features and four derived features, with a target variable consisting of 8 product classes. The sample size was chosen to satisfy two standard statistical constraints. First, a common rule of thumb in supervised tabular classification requires at least 10--20 observations per predictor variable \cite{ref52}. With 20 features, this heuristic would require a minimum of approximately 400 observations; 5{,}000 observations exceed this minimum by more than an order of magnitude. Second, the multi-class extension of the events-per-variable rule \cite{ref53} recommends at least 10--20 events per predictor variable per class to stabilize class-specific coefficients. Assuming near-balanced classes, 5{,}000 observations yield approximately 625 samples per class on average and therefore approximately 31 samples per feature per class, which is within the recommended range. Synthetic data generation permits arbitrary sample sizes because observations are produced by a data-generating process rather than collected from a finite population; the chosen size balances statistical coverage of the 8-class feature space with the computational cost of repeated cross-validation and interaction-detection passes \cite{ref52}, \cite{ref83}.

				\subsubsection{Feature Distributions and Parameters}
					The dataset contains 16 input features whose marginal distributions are grounded in published Philippine statistical sources. Continuous financial variables such as monthly income, monthly expenses, and existing savings follow non-uniform, right-skewed distributions consistent with what the literature on household economics reports for developing economies. Monthly income is sampled from a lognormal distribution whose location and scale parameters depend on employment status, age, and location, reflecting the well-documented empirical observation that income is approximately lognormally distributed across populations \cite{ref54}. Expenses are derived from income using a Beta-distributed expense ratio clipped to the interval $[0.40, 0.98]$, because the ratio of expenses to income is bounded between 0 and 1 and the Beta distribution is the standard choice for random variables on a bounded interval \cite{ref54}. Existing savings are generated as a function of monthly disposable income scaled by an exponentially distributed number of ``savings-months,'' consistent with the highly right-skewed buffer-stock savings distributions documented in BSP surveys \cite{ref35}. Age is modeled as a four-component Gaussian mixture clipped to $[18, 65]$, reflecting the Philippine demographic pyramid \cite{ref39}. Categorical features (employment status, education, location type, primary e-wallet, savings goal, risk tolerance, investment horizon) and binary features (\emph{has\_bank\_account}, \emph{has\_ewallet}, \emph{receives\_remittance}) are sampled from conditional categorical distributions whose probabilities are drawn from official surveys.

					The 16-feature schema includes three variables (\emph{education}, \emph{primary\_ewallet}, and \emph{receives\_remittance}) added in response to defense-panel feedback, who recommended (i)~introducing an additional discriminator to separate the previously confused \emph{maya\_savings} and \emph{gcash\_gsave} classes and (ii)~drawing on additional PSA survey items relevant to Filipino financial behavior. The \emph{primary\_ewallet} feature directly addresses (i) by exposing the platform-attachment signal that the original feature set could not observe; \emph{education} (PSA 2020 Census educational-attainment distribution \cite{ref39} and PSA LFS \cite{ref41}) and \emph{receives\_remittance} (PSA FIES \cite{ref40} and BSP FIS \cite{ref35} reporting on remittance-receiving households) address (ii). The integration of these variables and their downstream effect on per-class performance is documented in the Results and Discussion chapter.

					Table~\ref{tab:features} documents the statistical distribution, parameter set, and ground-truth source for each of the 16 input features.
{\centering\small
					\begin{xltabular}{\textwidth}{|>{\raggedright\arraybackslash}p{0.14\textwidth}|X|X|X|>{\centering\arraybackslash}p{0.12\textwidth}|}
						\caption{Feature Distributions and Parameter Justifications for the Synthetic Dataset}\label{tab:features}\\
						\hline
						\textbf{Feature} & \textbf{Distribution} & \textbf{Parameters} & \textbf{Justification} & \textbf{Source} \\
						\hline
						\endfirsthead
						\caption[]{Feature Distributions and Parameter Justifications for the Synthetic Dataset (\emph{continued})}\\
						\hline
						\textbf{Feature} & \textbf{Distribution} & \textbf{Parameters} & \textbf{Justification} & \textbf{Source} \\
						\hline
						\endhead
						age & Gaussian mixture (4 components), clipped [18, 65] & Modes at 24, 30, 40, 53 with stds 3, 5, 7, 6 and mixture weights (0.30, 0.30, 0.25, 0.15) & Captures Philippine age pyramid with a dominant young-adult mode & \cite{ref39} \\
						\hline
						monthly\_\newline income & Lognormal, parameters vary by employment × age × location & Base means PHP 3{,}000--25{,}000 by employment; age multipliers 0.70--1.40; location multipliers 0.75--1.40 & Right-skewed income distribution matching FIES family income & \cite{ref40}, \cite{ref54} \\
						\hline
						monthly\_\newline expenses & Derived: income $\times$ Beta ratio & Beta($\alpha{=}5, \beta{=}2) \times 0.5 + 0.4$, clipped [0.40, 0.98], $+$0.03 per dependent & Bounded-ratio structure; aligns with FIES expenditure patterns & \cite{ref40}, \cite{ref55} \\
						\hline
						existing\_\newline savings & Exponential-scaled disposable income, clipped to $\geq 0$ & months $\sim$ Exp(scale $=$ 6), clipped [0, 60]; savings $=$ disposable $\times$ months $+ \mathcal{N}(0, 5000)$ & Right-skewed savings buffer matching BSP FIS findings & \cite{ref35}, \cite{ref66} \\
						\hline
						employment\_status & Conditional categorical (5 levels) on age, location & See generation script for age-by-location probability tables & Labor-force composition from PSA LFS and freelancing prevalence & \cite{ref41} \\
						\hline
						education & Conditional categorical (5 levels) on age, location, employment & Levels: Elementary, High School, Vocational, College, Graduate; weighted by location (Metro Manila skews higher) and employed/student status & PSA 2020 Census educational-attainment distribution and PSA LFS labor-by-education tables & \cite{ref39}, \cite{ref41} \\
						\hline
						num\_\newline dependents & Conditional categorical on age & Peaks at 2--3 for ages 35--49; declines for under-25 and 50+ & Matches PSA household size ($\sim$4.1) and life-cycle composition & \cite{ref39} \\
						\hline
						location\_\newline type & Categorical (3 levels) & Metro Manila 13\%, Urban 45\%, Rural 42\% & Aligns with PSA Census urban/rural classification & \cite{ref39} \\
						\hline
						digital\_\newline savviness & Conditional ordinal (1--5) on age, location & Young urban cohort skews high (4--5); older rural cohort skews low (1--2) & Matches Digital 2023 Philippines penetration by demographic & \cite{ref42} \\
						\hline
						has\_bank\_\newline account & Conditional Bernoulli on employment, income, location & Base 0.56; $+$0.15 employed; $-$0.25 unemployed; $\pm$0.15 by location; $\pm$0.15 by income & 56\% national account ownership (BSP FIS) and Findex correlates & \cite{ref35}, \cite{ref66} \\
						\hline
						has\_ewallet & Conditional Bernoulli on age, location, digital\_savviness & Base 0.55; $\pm$0.20 by age band; $\pm$0.15 by location; $\pm 0.08$ per unit of digital savviness from 3 & $>$50\% e-wallet penetration (BSP FIS 2021) & \cite{ref35}, \cite{ref42} \\
						\hline
						primary\_ewallet & Conditional categorical (4 levels: GCash, Maya, Both, Neither) on \emph{has\_ewallet}, age, location & GCash leading at $\sim$32\% population share; Maya $\sim$16\%; Both $\sim$12\%; Neither $\sim$40\% (mirrors non-ewallet adopters) & GCash and Maya leading active-user shares as documented in BSP FIS and provider disclosures & \cite{ref9}, \cite{ref35}, \cite{ref42} \\
						\hline
						receives\_remittance & Conditional Bernoulli on location and income tier & Base 0.20; $+$0.10 rural; $+$0.05 lower-income; $-$0.05 high-income & PSA FIES remittance-receiving households ($\sim$22--25\%) and BSP FIS remittance-channel reporting & \cite{ref35}, \cite{ref40} \\
						\hline
						savings\_goal & Conditional categorical (7 levels) on age, employment & Students favor Education; young adults favor Emergency Fund; mid-career favor Home/Property; older favor Retirement & Life-cycle goal evolution aligned with LCH & \cite{ref65} \\
						\hline
						risk\_tolerance & Conditional categorical (3 levels) on age, income, dependents & Aggressive peaks at young/high-income/low-dependents; Conservative peaks at older/low-income/high-dependents & Prospect theory; risk aversion rises with age and obligations & \cite{ref62} \\
						\hline
						investment\_horizon & Conditional categorical (3 levels) on savings\_goal & Retirement $\to$ Long-term; Emergency/Travel $\to$ Short-term; Home/Education $\to$ Medium-term & Horizon dictated by goal per intertemporal choice theory & \cite{ref64}, \cite{ref84} \\
					\end{xltabular}\par}

					The lognormal distribution is adopted for income because income is strictly positive, multiplicative in its determinants (education, experience, location premiums), and right-skewed across populations; these are precisely the conditions under which Limpert, Stahel, and Abbt \cite{ref54} argue the lognormal is a more appropriate model than the normal. The Beta distribution is adopted for the expense ratio because it is the canonical parametric family for random variables on a bounded interval $[0,1]$ and admits any desired mean and skewness via its $\alpha$ and $\beta$ parameters, making it the standard choice for modeling proportions and ratios in applied statistics \cite{ref54}.

				\subsubsection{Conditional Dependency Structure}
					A central concern in synthetic tabular data generation is that independently sampled features break inter-feature correlations present in real populations \cite{ref29}. To address this, features in the FinXplain dataset are generated in topological order according to a directed dependency structure. Specifically: \emph{location\_type} and \emph{age} are sampled first; \emph{employment\_status}, \emph{num\_dependents}, \emph{digital\_savviness}, \emph{risk\_tolerance}, and \emph{savings\_goal} are then sampled conditional on \emph{age} and/or \emph{location\_type}; \emph{monthly\_income} is sampled conditional on \emph{employment\_status}, \emph{age}, and \emph{location\_type}; \emph{monthly\_expenses} is derived from \emph{monthly\_income} and \emph{num\_dependents} via a Beta-distributed expense ratio; \emph{existing\_savings} is derived from \emph{monthly\_income} and \emph{monthly\_expenses}; \emph{has\_bank\_account} is conditioned on \emph{employment\_status}, \emph{monthly\_income}, and \emph{location\_type}; \emph{has\_ewallet} is conditioned on \emph{age}, \emph{location\_type}, and \emph{digital\_savviness}; \emph{primary\_ewallet} is conditioned on \emph{has\_ewallet}, \emph{age}, and \emph{location\_type} (and is forced to ``Neither'' whenever \emph{has\_ewallet} $=0$); \emph{education} is conditioned on \emph{age}, \emph{location\_type}, and \emph{employment\_status}; and \emph{receives\_remittance} is conditioned on \emph{location\_type} and the \emph{monthly\_income} tier. This dependency Directed Acyclic Graph (DAG) encodes the empirically documented relationships between demographic, geographic, and economic variables reported in Philippine statistical sources \cite{ref35}, \cite{ref40}, \cite{ref66}.

					This hand-encoded dependency structure is a simplification of the fully learned correlation structure produced by generative models such as CTGAN \cite{ref34} or the Synthetic Data Vault \cite{ref29}, but it has the advantage of being auditable: every conditional probability table in the generator is directly traceable to a named statistical source. Residual correlations that are not explicitly encoded in the DAG are partially induced through the rule-based label-generation logic described in the next subsection, because multiple features jointly enter the product-assignment score functions and thus correlate with the assigned label even when they are pairwise independent in the feature matrix.

				\subsubsection{Derived Features}
					Four derived features are computed from the 16 input features during preprocessing:

					\begin{itemize}
					    \item \textbf{\emph{disposable\_income} $=$ \emph{monthly\_income} $-$ \emph{monthly\_expenses}.} This is the absolute monthly surplus available for savings or investment and is the direct determinant of savings in Keynesian consumption theory and the life-cycle hypothesis of Modigliani and Brumberg \cite{ref65}.
					    \item \textbf{\emph{expense\_ratio} $=$ \emph{monthly\_expenses} $/$ (\emph{monthly\_income} $+ 1$).} This is the fraction of income consumed. The complement, $1 -$ \emph{expense\_ratio}, is the savings rate, a central variable in household consumption--saving models \cite{ref55}.
					    \item \textbf{\emph{savings\_to\_income\_ratio} $=$ \emph{existing\_savings} $/$ (\emph{monthly\_income} $+ 1$).} This expresses the existing savings buffer in months-of-income units and is a standard measure of buffer-stock adequacy in household-finance literature \cite{ref55}.
					    \item \textbf{\emph{income\_per\_dependent} $=$ \emph{monthly\_income} $/$ (\emph{num\_dependents} $+ 1$).} This is a per-capita income proxy adjusted for household composition, following the equivalence-scale logic used in welfare economics \cite{ref55}.
					\end{itemize}

					The $+1$ offset in the three ratio features prevents division by zero. These derived features are consumed directly by the EBM model alongside the 16 input features and by the rule-based label generator; they are not used as secondary target variables.

				\subsubsection{Label Generation Logic}
					\label{sec:label-gen}
					Product labels are not assigned by uniform random sampling or by a single deterministic threshold. They are assigned by a multi-factor additive scoring system in which each candidate product accumulates a score across eleven scoring dimensions (platform attachment via \emph{primary\_ewallet}, digital preference, income level, existing savings, risk tolerance, investment horizon, savings goal, location, age, education, and remittance receipt), each grounded in the product's official eligibility criteria and in established financial-suitability principles. A platform-attachment bonus of up to $+3$ is added to the products of the user's declared primary e-wallet (\emph{GCash} $\to+3$ \emph{gcash\_gsave}; \emph{Maya} $\to+3$ \emph{maya\_savings}; \emph{Both} $\to+1$ to all GCash/Maya digital savings; \emph{Neither} $\to+2$ \emph{bpi\_savings}, with a smaller bonus to other BPI products), reflecting brand-driven adoption patterns reported in Philippine fintech research \cite{ref9} \emph{using a feature the user actually inputs} rather than a hidden randomly assigned preference. This change replaces the earlier hidden-random platform-preference bonus, after defense-panel feedback that the previously confused \emph{maya\_savings} and \emph{gcash\_gsave} classes required an additional observable discriminator. A stochastic noise factor of 8\% replaces the top-scoring product with a random choice among the top-3 scoring products, preventing unrealistically deterministic labels and providing mild regularization to downstream classifiers \cite{ref46}.

					The decision rules can be expressed as conditional IF-THEN statements derived directly from the generator code. Representative rules are summarized below, with product-specific justifications.
					\begin{itemize}
					    \item \textbf{R1 -- GSave (GCash):} \textsc{If} \emph{existing\_savings} $<$ PHP 5{,}000 \textsc{and} (\emph{digital\_savviness} $\geq 3$ \textsc{or} \emph{has\_ewallet} $= 1$) \textsc{then} favor \emph{gcash\_gsave}. GSave has no minimum balance, is accessible entirely through the GCash application, and is PDIC-insured, making it suitable for low-savings, digitally engaged users \cite{ref35}. Here ``low \emph{existing\_savings}'' corresponds to the below-PHP-5{,}000 tier in the generator's additive scoring function, which is the tier at which \emph{gcash\_gsave} receives its largest savings-dimension bonus.
					    \item \textbf{R2 -- GInvest (GCash):} \textsc{If} \emph{risk\_tolerance} $=$ Aggressive \textsc{and} \emph{investment\_horizon} $=$ Long-term \textsc{then} favor \emph{gcash\_ginvest}. GInvest UITFs are market-linked and appropriate only for users who tolerate market risk and have a long holding horizon \cite{ref62}, \cite{ref84}.
					    \item \textbf{R3 -- Maya Savings:} \textsc{If} \emph{monthly\_income} is low to moderate \textsc{and} \emph{primary\_ewallet} $\in$ \{Maya, Both\} \textsc{and} \emph{digital\_savviness} $\geq 3$ \textsc{then} favor \emph{maya\_savings}. Maya's base rate is several multiples of traditional bank savings rates, and its digital-first delivery suits moderate-income digitally engaged users; the \emph{primary\_ewallet} requirement encodes the user's platform attachment as the discriminator that separates \emph{maya\_savings} from the otherwise feature-equivalent \emph{gcash\_gsave} \cite{ref9}, \cite{ref35}.
					    \item \textbf{R4 -- Maya Personal Goals:} \textsc{If} \emph{savings\_goal} $\in$ \{Education, Home/Property, Travel\} \textsc{and} \emph{investment\_horizon} $=$ Medium-term \textsc{then} favor \emph{maya\_personal\_goals}. Goal-based savings products operationalize mental-accounting and commitment-device findings \cite{ref63}.
					    \item \textbf{R5 -- Maya Time Deposit Plus:} \textsc{If} \emph{risk\_tolerance} $\neq$ Aggressive \textsc{and} PHP 30{,}000 $\leq$ \emph{monthly\_income} $\leq$ PHP 60{,}000 \textsc{and} \emph{investment\_horizon} $=$ Medium-term \textsc{then} favor \emph{maya\_time\_deposit}. A short-lock digital time deposit with above-market fixed rates is appropriate for moderate-income conservative users with medium horizons \cite{ref35}.
					    \item \textbf{R6 -- BPI Regular Savings:} \textsc{If} \emph{has\_bank\_account} $= 1$ \textsc{and} \emph{risk\_tolerance} $=$ Conservative \textsc{and} \emph{digital\_savviness} $\leq 3$ \textsc{then} favor \emph{bpi\_savings}. The product is a legacy bank-branch/ATM account appropriate for conservative users with existing bank relationships \cite{ref37b}.
					    \item \textbf{R7 -- BPI~\#SaveUp:} \textsc{If} \emph{has\_bank\_account} $= 1$ \textsc{and} \emph{digital\_savviness} $\geq 3$ \textsc{and} \emph{savings\_goal} $\neq$ Retirement \textsc{then} favor \emph{bpi\_save\_up}. A digital goal-based auto-save within the BPI ecosystem fits young, formally banked, digitally engaged users \cite{ref37b}, \cite{ref63}.
					    \item \textbf{R8 -- BPI Plan Ahead Time Deposit:} \textsc{If} \emph{existing\_savings} $\geq$ PHP 50{,}000 \textsc{and} \emph{risk\_tolerance} $=$ Conservative \textsc{and} \emph{investment\_horizon} $=$ Long-term \textsc{then} favor \emph{bpi\_time\_deposit}. Justification: The 5-year fixed-term deposit has a PHP 50{,}000 minimum and is appropriate only for users with substantial existing savings and long-horizon preferences \cite{ref37b}, \cite{ref64}.
					    \item \textbf{R9 -- Education adjustment:} \textsc{If} \emph{education} $\in$ \{College, Graduate\} \textsc{then} add bonuses to \emph{gcash\_ginvest}, \emph{bpi\_time\_deposit}, and \emph{maya\_time\_deposit} (more complex products with explicit terms); \textsc{if} \emph{education} $\in$ \{Elementary, High School\} \textsc{then} penalize \emph{gcash\_ginvest} and shift weight toward simple no-minimum digital savings (\emph{gcash\_gsave}, \emph{maya\_savings}). This reflects the financial-literacy-by-attainment relationship documented by Lusardi and Mitchell \cite{ref43} and the PSA LFS attainment-by-occupation tabulations \cite{ref41}.
					    \item \textbf{R10 -- Remittance receipt:} \textsc{If} \emph{receives\_remittance} $= 1$ \textsc{then} add bonuses to \emph{maya\_personal\_goals}, \emph{bpi\_save\_up}, \emph{bpi\_time\_deposit}, and \emph{maya\_time\_deposit}. Remittance-recipient households in PSA FIES \cite{ref40} display systematically higher household saving rates and a documented preference for goal-based and time-locked vehicles that discipline irregular inflows \cite{ref35}, \cite{ref63}.
					\end{itemize}

					Rule-based label generation for supervised classification has strong precedent in the financial machine-learning literature. West \cite{ref56} uses expert-rule-derived credit-scoring labels as a benchmark for neural-network credit models; Thomas \cite{ref57} reviews the long history of rule-based and scorecard-based financial decision systems; and content-based recommender-systems research commonly uses expert rules to create ground-truth labels when no behavioral data is available \cite{ref13}. As Assefa et al. \cite{ref46} caution, rule-based labels are a simplification of the full complexity of real bank eligibility decisions and should be explicitly documented as such; they are used here precisely because no real labels exist for Filipino fintech consumers.

					A natural question is whether training a classifier on rule-generated labels contributes anything beyond a re-execution of the rules themselves. Four properties of the FinXplain pipeline establish the ML model as a non-trivial extension of the rule base rather than a re-encoding of it. First, the additive scoring system assigns \emph{every} product a continuous score rather than a single class, and the trained EBM preserves this structure by producing a probability distribution over all eight products; this distribution powers the top-3 product listing in the prototype interface, which the discrete rules cannot produce without further post-processing. Second, the 8\% stochastic top-3 substitution makes labels non-deterministic at the decision boundary: at least 8\% of labels disagree with the \textsc{argmax} of the scoring function, so a rule-direct classifier cannot recover the observed labels exactly and the ML model must learn the \emph{conditional} label distribution rather than the \textsc{argmax} rule. Third, the EBM's univariate shape functions \cite{ref48}, \cite{ref49} provide a continuous, visualizable representation of how each feature contributes across its entire domain, whereas the discrete IF--THEN rules only expose the endpoints of each condition (e.g.\ ``\emph{monthly\_income} $\geq$ PHP 30{,}000''); the shape functions can thus surface gradient effects that the rules are deliberately too coarse to express. Fourth, the EBM's pairwise interaction terms \cite{ref15} can learn residual feature--feature correlations that the additive scoring system does not explicitly encode, enabling the model to generalize to novel feature combinations that fall between the hand-specified rule cases. These four properties, taken together, position the ML model as a continuous, probabilistic, and generalizable surrogate for a deliberately coarse rule-based oracle rather than as a memorizer of the oracle itself.

					The resulting class distribution is emergent from the interaction of the feature distributions and the scoring logic rather than manually specified. Stratified 80/20 splitting is used so that each of the eight product classes is represented proportionally in the training and test sets, and stratified five-fold cross-validation is used for generalization estimation \cite{ref59}. When class imbalance is observed in the generated labels, the weighted-averaging evaluation strategy described in the Results chapter accounts for imbalance at evaluation time; this is the approach recommended for imbalanced multiclass problems \cite{ref58}.

				\subsubsection{Feature-to-Outcome Relationships}
					Table 3 summarizes the expected direction of effect of each feature on product recommendation, grounded in financial theory.

					\vspace{2ex}
					{\centering\small
					\begin{xltabular}{\textwidth}{|>{\raggedright\arraybackslash}p{0.17\textwidth}|>{\centering\arraybackslash}p{0.06\textwidth}|X|X|X|>{\centering\arraybackslash}p{0.11\textwidth}|}
						\caption{Expected Feature-to-Outcome Relationships Based on Financial Theory}\label{tab:feat-outcome}\\
						\hline
						\textbf{Feature} & \textbf{Dir.} & \textbf{Products Favored} & \textbf{Products Disfavored} & \textbf{Theoretical Basis} & \textbf{Source} \\
						\hline
						\endfirsthead
						\caption[]{Expected Feature-to-Outcome Relationships Based on Financial Theory (\emph{continued})}\\
						\hline
						\textbf{Feature} & \textbf{Dir.} & \textbf{Products Favored} & \textbf{Products Disfavored} & \textbf{Theoretical Basis} & \textbf{Source} \\
						\hline
						\endhead
						risk\_tolerance (Aggressive) & $+$ & GInvest & BPI Savings, BPI TD & Prospect theory; risk preferences & \cite{ref62} \\
						\hline
						investment\_\newline horizon (Long) & $+$ & BPI TD, GInvest & GSave, Maya Savings & Time diversification and horizon effects & \cite{ref84} \\
						\hline
						savings\_goal (goal-named) & $+$ & Maya Personal Goals, BPI~\#SaveUp & Plain savings & Goal-based investing; commitment devices & \cite{ref63}, \cite{ref85} \\
						\hline
						monthly\_\newline income (high) & $+$ & BPI TD, GInvest & GSave, Maya Savings & Permanent income hypothesis; risk capacity & \cite{ref64} \\
						\hline
						age (older) & $+$ & BPI products, TDs & GInvest, digital savings & Life-cycle hypothesis & \cite{ref65} \\
						\hline
						digital\_\newline savviness (high) & $+$ & GCash, Maya products & BPI Regular Savings & Digital financial-literacy linkage & \cite{ref66} \\
						\hline
						has\_ewallet & $+$ & GCash, Maya products & BPI Regular Savings, BPI TD & E-wallet ownership as digital channel & \cite{ref35}, \cite{ref66} \\
						\hline
						has\_bank\_\newline account & $+$ & BPI products & GSave, Maya Savings (as sole option) & Bank account as prerequisite for BPI & \cite{ref35}, \cite{ref66} \\
						\hline
						num\_\newline dependents (high) & $+$ & Liquid, low-risk savings & GInvest, long-lock TD & Household composition and risk aversion & \cite{ref55} \\
						\hline
						location\_type (rural) & $+$ & GCash, Maya (digital-first) & BPI branch-based & Geographic access barriers & \cite{ref35} \\
						\hline
						employment\_\newline status (employed) & $+$ & BPI products (payroll) & --- & Income stability; product eligibility & \cite{ref64}, \cite{ref41} \\
						\hline
						primary\_\newline ewallet (GCash) & $+$ & GSave, GInvest & Maya products & Platform attachment as observable adoption signal & \cite{ref9}, \cite{ref35} \\
						\hline
						primary\_\newline ewallet (Maya) & $+$ & Maya Savings, Personal Goals, TD Plus & GCash products & Platform attachment; resolves Maya/GCash overlap & \cite{ref9}, \cite{ref35} \\
						\hline
						primary\_ewallet (Neither) & $+$ & BPI products & GCash, Maya digital products & Non-adopters fall back to traditional banking & \cite{ref35}, \cite{ref66} \\
						\hline
						education (College/\newline Graduate) & $+$ & GInvest, Time Deposits & Plain digital savings & Financial-literacy gradient; complex-product engagement & \cite{ref43}, \cite{ref41} \\
						\hline
						receives\_\newline remittance & $+$ & Maya Personal Goals, BPI~\#SaveUp, Time Deposits & --- & Goal-based and disciplined-saving preference of remittance recipients & \cite{ref40}, \cite{ref63} \\
					\end{xltabular}\par}
					\vspace{4ex}

					The EBM model's shape functions, which are a core feature of the Explainable Boosting Machine architecture \cite{ref48}, \cite{ref49}, allow direct visualization of how each feature contributes to the recommendation of each product class. Inspecting these shape functions is the primary mechanism for confirming or refuting the theoretical directions summarized above; a misaligned shape function is a flag that either the label-generation rules or the synthetic-data parameters may not match financial theory.

					SHAP values \cite{ref18}, computed on the test set during offline analysis of the XGBoost benchmark, provide complementary empirical evidence of feature-importance ranking. The top features identified by SHAP analysis of the benchmark (most prominently \emph{primary\_ewallet}, monthly income, risk tolerance, existing savings, and digital savviness) are consistent with the theoretically expected directions documented in Table~\ref{tab:feat-outcome}; the prominence of \emph{primary\_ewallet} in particular is the empirical confirmation that the panel-recommended platform-attachment discriminator is doing the work it was introduced to do. This consistency across the interpretable EBM and the black-box benchmark is a face-validity check on the synthetic dataset: it indicates that the learned relationships track documented financial behavior rather than artifacts of the data-generating process.

				\subsubsection{Data Reliability and Limitations}
					The entire generation pipeline is deterministic given the fixed random seed $42$, so the dataset is bit-for-bit reproducible from the script \cite{ref60}. The train/test split uses \path{stratify=y} so that all 8 classes are represented in the training and test partitions, and stratified five-fold cross-validation is used on the full dataset; these are the standard techniques for reliable evaluation under potential class imbalance \cite{ref59}.

					Each feature distribution is parameterized against a named Philippine or international source (PSA Census \cite{ref39}, PSA FIES \cite{ref40}, PSA Labor Force Survey \cite{ref41}, BSP Financial Inclusion Survey \cite{ref35}, Digital 2023 Philippines \cite{ref42}, Global Findex \cite{ref66}). This grounds the marginal distributions of the synthetic population in empirically documented population-level statistics, although the mapping from population-level statistics to individual-level synthetic records is itself an assumption that cannot be fully validated without real micro-data.

					The synthetic dataset has not been validated against a held-out sample of real Filipino fintech consumers, and results on synthetic data cannot be assumed to generalize to a production environment without subsequent validation \cite{ref46}, \cite{ref29}. The ``model card'' best practice of documenting this explicitly \cite{ref61} is adopted here: the trained model is positioned as a research prototype, not a deployable recommender.

					The following limitations are acknowledged by the researcher in the study:

					\begin{itemize}
					    \item Rule-based label generation simplifies real eligibility decisions and may not capture the full complexity of bank suitability assessments \cite{ref46}.
					    \item Feature generation is conditionally structured but not jointly distributed; some higher-order correlations present in real data may be under-represented \cite{ref29}, \cite{ref34}.
					    \item No empirical validation against real Philippine bank customer data has been performed.
					    \item When class imbalance is observed in the generated labels, the weighted-averaging strategy mitigates but does not eliminate the evaluation bias it introduces \cite{ref58}.
					    \item Interest rates, minimum balances, and product features used in the label rules are point-in-time values (April 2026) and are subject to change by the providers.
					\end{itemize}

					Label-rule simplicity is partially mitigated by the eight-dimension additive scoring and the 8\% noise injection, which soften decision boundaries. The absence of learned joint dependencies is partially compensated by the explicit conditional DAG and by the EBM's pairwise interaction terms, which can learn residual correlations. The absence of real-data validation is explicitly flagged as future work; the BSP Financial Consumer Protection Act \cite{ref44} anticipates that such validation would be a precondition for deployment. Class imbalance is handled through stratified splitting and weighted-averaged metrics \cite{ref58}, \cite{ref78}. Drift in product rates is mitigated by separating product metadata (\emph{product\_info.json}) from recommendation logic, so that rate updates do not require retraining the model.

			\subsection{XAI Methodologies}
				\subsubsection{Explainable Boosting Machine (EBM)}
					The Explainable Boosting Machine is a glass-box machine-learning model based on Generalized Additive Models with pairwise interactions (GA$^2$M), developed at Microsoft Research \cite{ref48}, \cite{ref49}, \cite{ref15}. An EBM approximates the link-transformed conditional expectation as
					\begin{equation*}
					    g(E[y \mid x]) = \beta_0 + \sum_{i} f_i(x_i) + \sum_{(i,j) \in \mathcal{P}} f_{ij}(x_i, x_j),
					\end{equation*}
					where $g$ is the softmax link for multiclass classification, each univariate shape function $f_i$ is learned by cyclic gradient boosting of shallow trees on residuals, and each pairwise interaction function $f_{ij}$ is learned only for a small set $\mathcal{P}$ of feature pairs selected automatically via FAST interaction detection \cite{ref49}. Because the prediction is an additive sum of independently inspectable shape functions, each feature's contribution to any single prediction can be read off directly from the model without any post-hoc approximation \cite{ref50}, \cite{ref15}.

					The EBM is trained using InterpretML's \emph{ExplainableBoostingClassifier} with the following hyperparameters (verified in \emph{notebooks/train\_ebm.py}): \emph{max\_bins}$=256$, \emph{interactions}$=10$, \emph{outer\_bags}$=8$, \emph{inner\_bags}$=0$, \emph{learning\_rate}$=0.01$, \emph{max\_rounds}$=5000$, \emph{min\_samples\_leaf}$=2$, \emph{random\_state}$=42$, and an explicit \emph{feature\_names} list of 20 entries (16 input features plus 4 derived features). The remaining arguments use InterpretML defaults \cite{ref15}.

					The EBM was chosen as the primary production model for four reasons. First, it provides native interpretability without post-hoc approximation, avoiding the known unfaithfulness of post-hoc explanation methods \cite{ref51}. Second, it is appropriate for financial decision support where transparency is legally required: the Philippine Financial Consumer Protection Act (RA~11765) \cite{ref44} and BSP consumer-protection circulars mandate that financial product disclosures be intelligible to the consumer, and the Data Privacy Act (RA~10173) \cite{ref45} requires disclosure of the logic behind automated decisions. Third, on tabular data EBMs (and the closely related intelligible-GAM models that share the same additive structure) have been demonstrated to achieve accuracy competitive with gradient-boosted ensembles while remaining interpretable \cite{ref15}, \cite{ref48}, \cite{ref49}, \cite{ref50}. Fourth, Rudin \cite{ref51} argues that for high-stakes decisions on tabular data, practitioners should prefer inherently interpretable models over black-box models with post-hoc explanations; the EBM's glass-box architecture is the direct embodiment of that argument.

				\subsubsection{XGBoost (Benchmark Model)}
					XGBoost (eXtreme Gradient Boosting) is a scalable regularized variant of gradient tree boosting \cite{ref16}, \cite{ref67}. For multi-class classification it minimizes an objective of the form
					\begin{equation*}
					    \mathcal{L}(\theta) = \sum_{i} \ell(\hat{y}_i, y_i) + \sum_{k} \Omega(f_k),
					\end{equation*}
					where $\ell$ is the multi-class softmax log-loss and $\Omega(f) = \gamma T + \tfrac{1}{2}\lambda \| w \|^2$ penalizes tree complexity via the number of leaves $T$ and the magnitude of leaf weights $w$. Trees are added sequentially, each fitting the negative gradient of the loss with respect to the current ensemble's predictions, and a second-order Taylor expansion is used to compute optimal split points and leaf weights \cite{ref16}, \cite{ref67}.

					The XGBoost benchmark is trained with \emph{n\_estimators}$=300$, \emph{max\_depth}$=6$, \emph{learning\_rate}$=0.1$, \emph{subsample}$=0.8$, \emph{colsample\_bytree}$=0.8$, \emph{min\_child\_weight}$=3$, \emph{gamma}$=0.1$, \emph{reg\_alpha}$=0.1$, \emph{reg\_lambda}$=1.0$, \emph{objective}$=$\emph{multi:softprob}, \emph{num\_class}$=8$, \emph{eval\_metric}$=$\emph{mlogloss}, and \emph{random\_state}$=42$. XGBoost requires explicit preprocessing in the benchmark pipeline: categorical features are integer-encoded via \emph{OrdinalEncoder} and numerical features are centered/scaled via \emph{StandardScaler} \cite{ref74}.

					XGBoost is used exclusively as a benchmark model in this study; it is \emph{not} in the production \emph{/recommend} endpoint. Its role is to validate the predictive quality of the EBM on the held-out test set under identical evaluation conditions \cite{ref68}. This design reflects two considerations. First, XGBoost is widely regarded as a state-of-the-art classifier for tabular data \cite{ref16}, so it provides a strong upper bound against which the interpretable EBM can be compared \cite{ref51}. Second, XGBoost is well documented, widely adopted, and reproducible, which matches the methodological-comparison role it plays here. For a prototype whose stated purpose is explanation quality under regulatory transparency constraints \cite{ref44}, deploying the model whose own internals are the explanation is the more defensible choice.

				\subsubsection{SHAP Explanations}
					SHAP (SHapley Additive exPlanations) assigns each feature a contribution value derived from Shapley values \cite{ref70}, which allocate the output of a cooperative game among its players according to each player's marginal contribution across all possible coalitions. Lundberg and Lee \cite{ref18} unify a family of additive feature-attribution methods under the SHAP framework, in which any local explanation of a prediction $f(x)$ takes the additive form
					\begin{equation*}
					    f(x) \approx \phi_0 + \sum_{i=1}^{M} \phi_i,
					\end{equation*}
					where $\phi_0$ is the expected prediction and $\phi_i$ is the SHAP value for feature $i$. SHAP values are the unique feature attributions that simultaneously satisfy local accuracy, missingness, and consistency \cite{ref18}. For tree-based models, \emph{TreeSHAP} computes exact SHAP values in polynomial time by exploiting the tree structure \cite{ref69}.

					In the FinXplain codebase, SHAP is used in the offline benchmark script with \emph{shap.TreeExplainer} applied to the XGBoost model on the first 100 test samples. Per-class mean absolute SHAP values are computed to rank feature importance, and the top-5 features per class are printed. The production API does \emph{not} call SHAP at inference time; it uses the EBM's exact additive \emph{explain\_local()} output to populate the \emph{key\_factors} array that powers the horizontal bar chart on the \emph{/results} page. Each key-factor entry records the feature name, a human-readable label, the user's actual input value, the raw contribution, a direction tag (\emph{supports} if $\phi_i > 0$, \emph{opposes} if $\phi_i < 0$), and an impact tier (\emph{high} if $|\phi_i| \geq 1.5$, \emph{medium} if $|\phi_i| \geq 0.5$, \emph{low} otherwise). This per-feature breakdown gives users a quantitative, ranked view of what drove their recommendation.

				\subsubsection{LIME Explanations}
					LIME (Local Interpretable Model-agnostic Explanations) explains an individual prediction by fitting a sparse linear surrogate model in the neighborhood of the input instance. Ribeiro, Singh, and Guestrin \cite{ref19} formalize LIME as
					\begin{equation*}
					    \xi(x) = \arg\min_{g \in \mathcal{G}}\; \mathcal{L}(f, g, \pi_x) + \Omega(g),
					\end{equation*}
					where $f$ is the black-box model, $g$ is a surrogate drawn from an interpretable family $\mathcal{G}$, $\pi_x$ is a locality kernel around the instance $x$, $\mathcal{L}$ measures the surrogate's local fidelity to $f$, and $\Omega(g)$ penalizes surrogate complexity. LIME is model-agnostic: it only needs query access to $f$. Garreau and von Luxburg \cite{ref71} provide a theoretical analysis of LIME's behavior and document its known sensitivity to the choice of neighborhood kernel and perturbation distribution.

					In this study, LIME is implemented via \emph{LimeTabularExplainer} in the offline benchmark script, with \emph{num\_features}$=10$, \emph{top\_labels}$=3$, the training set used to fit the perturbation distribution, explicit categorical feature indices, and \emph{random\_state}$=42$. It is applied to the XGBoost model for methodological comparison against SHAP and against the EBM's native explanations. In production, the human-readable plain-English \emph{top\_reasons} and \emph{against\_reasons} displayed on the results page are generated by a deterministic template-based system over EBM contributions rather than by live LIME queries; LIME's role is therefore strictly analytic in this codebase.

				\subsubsection{Counterfactual Explanations}
					Counterfactual explanations answer the question ``what minimal change to the input would produce a different model output?'' Wachter, Mittelstadt, and Russell \cite{ref20} formalize the counterfactual as
					\begin{equation*}
					    \arg\min_{x'}\;\max_\lambda\; \lambda (f(x') - y')^2 + d(x, x'),
					\end{equation*}
					where $x'$ is the counterfactual instance, $y'$ is the desired (different) output, and $d(x, x')$ is a distance metric that penalizes large departures from the original input. Karimi, Sch\"olkopf, and Valera \cite{ref72} extend this to algorithmic recourse, distinguishing between contrastive explanations and actionable interventions, and emphasize that only interventions on features the user can actually control constitute recourse. Mothilal, Sharma, and Tan \cite{ref73} introduce diversity as a second desideratum alongside proximity, so that multiple qualitatively different counterfactuals are offered rather than small perturbations of a single one.

					The FinXplain counterfactual generator is a custom implementation in \emph{backend/explainability.py}. It does not call an external library such as DiCE \cite{ref73}. Features are partitioned into an actionable set---\emph{existing\_savings}, \emph{monthly\_income}, \emph{monthly\_expenses}, \emph{risk\_tolerance}, \emph{investment\_horizon}, \emph{savings\_goal}---and a fixed set---\emph{age}, \emph{location\_type}, \emph{employment\_status}, \emph{has\_bank\_account}, \emph{has\_ewallet}, \emph{num\_dependents}, \emph{digital\_savviness}---following the actionability principle of Ustun et al.\ \cite{ref86}. The fixed set captures features that are either demographic/structural (\emph{age}, \emph{location\_type}, \emph{num\_dependents}), tied to external institutional relationships (\emph{employment\_status}, \emph{has\_bank\_account}, \emph{has\_ewallet}), or best interpreted as a slowly-changing user attribute (\emph{digital\_savviness}); within the short-to-medium financial-planning horizon addressed by the prototype, these features are treated as immutable so that every generated counterfactual corresponds to guidance the user can realistically act on rather than to hypothetical changes that would require relocating, altering household composition, or changing employment status. For each actionable numerical feature, the generator applies multiplicative perturbations $\{1.25, 1.5, 2.0, 2.5, 3.0, 0.75, 0.5\}$ (or absolute targets $\{10{,}000, 25{,}000, 50{,}000, 100{,}000, 200{,}000\}$ when the current value is zero) and records the first perturbation that flips the predicted class. For each actionable categorical feature, the generator tests every alternative value in the feature's domain and records the first flip. At most one flip per feature is recorded, and at most three counterfactuals are returned in total, enforcing diversity across features \cite{ref73} and proximity within features \cite{ref20}. For each counterfactual, the alternative product is identified by re-predicting with the modified feature vector. The user sees the original value, the suggested value, the alternative product name, the alternative provider, and a plain-language ``if you changed X to Y, product Z would be recommended'' sentence, transforming the prototype from a passive recommender into an active financial-planning tool \cite{ref72}.

			\subsection{Evaluation Procedures}
				The evaluation procedure focuses exclusively on model performance on the held-out test set. Explainability-quality, user-study, and functional-validation evaluations are out of scope for this study.

				\subsubsection{Model Performance Evaluation Overview}
					Evaluation is conducted on the held-out test set produced by the stratified 80/20 split described above: 4{,}000 training observations, 1{,}000 test observations, stratified by the 8-class target, with \path{random_state=42}. Both the EBM and the XGBoost benchmark are evaluated on the same test set under identical conditions. Five metrics are reported: accuracy, weighted precision, weighted recall, weighted F1 score, and weighted one-vs-rest AUC-ROC. All metrics are computed using scikit-learn \cite{ref74}. Weighted averaging is used for all multiclass metrics to account for potential class imbalance in the synthetic product-label distribution; this is the strategy recommended by Sokolova and Lapalme \cite{ref78} for multiclass settings in which class sizes may differ.

				\subsubsection{Accuracy}
					Accuracy is the fraction of test observations whose predicted label matches the true label. It is computed via \emph{sklearn.metrics.accuracy\_score}. Accuracy answers: of all recommendations the model makes, what fraction are correct? It has well-known limitations under class imbalance, because it weights all classes equally regardless of prevalence \cite{ref78}, \cite{ref58}; for this reason it is reported together with precision, recall, F1, and AUC-ROC rather than in isolation. Accuracy is the most commonly reported metric in classification benchmarks \cite{ref87} and is included here for comparability with prior work.

				\subsubsection{Precision (Weighted)}
					For each class $c$, precision is the fraction of instances predicted as $c$ that actually belong to $c$:
					\begin{equation*}
					    \mathrm{Precision}_c = \frac{TP_c}{TP_c + FP_c}, \quad \mathrm{Precision}_{\mathrm{weighted}} = \sum_c \frac{n_c}{N}\,\mathrm{Precision}_c.
					\end{equation*}
					It is computed via \path{precision_score(y_test, y_pred, average='weighted')}. Weighted precision measures: when the model recommends a particular product, how often is that recommendation correct, averaged with weights proportional to class support \cite{ref79}, \cite{ref78}. High precision reduces the risk of mismatched recommendations, which is directly relevant in a financial-product setting.

				\subsubsection{Recall (Weighted)}
					Weighted recall is defined symmetrically:
					\begin{equation*}
					    \mathrm{Recall}_c = \frac{TP_c}{TP_c + FN_c}, \quad \mathrm{Recall}_{\mathrm{weighted}} = \sum_c \frac{n_c}{N}\,\mathrm{Recall}_c.
					\end{equation*}
					It is computed via \path{recall_score(y_test, y_pred, average='weighted')}. Recall answers: of all users who should have received product $c$, what proportion does the model correctly identify \cite{ref79}, \cite{ref78}? High recall ensures that no suitable product is systematically missed.

				\subsubsection{F1 Score (Weighted)}
					The F1 score is the harmonic mean of precision and recall:
					\begin{equation*}
					    F1_c = \frac{2\,\mathrm{Precision}_c\,\mathrm{Recall}_c}{\mathrm{Precision}_c + \mathrm{Recall}_c}, \quad F1_{\mathrm{weighted}} = \sum_c \frac{n_c}{N}\,F1_c.
					\end{equation*}
					It is computed via \path{f1_score(y_test, y_pred, average='weighted')}. F1 is the primary metric in this study because it balances precision and recall and is robust to class imbalance \cite{ref78}, \cite{ref80}; the imbalanced-multiclass classification literature recommends F1 (particularly macro- or weighted-averaged F1) as the headline metric over raw accuracy when class sizes differ \cite{ref58}.

				\subsubsection{AUC-ROC (Weighted One-vs-Rest)}
					The AUC-ROC is the area under the Receiver Operating Characteristic curve, which plots the true positive rate against the false positive rate at every classification threshold. For the multiclass setting, the one-vs-rest (OvR) strategy computes an AUC for each class against the union of the remaining classes and takes a support-weighted average:
					\begin{equation*}
					    \mathrm{AUC}_{\mathrm{weighted}} = \sum_c \frac{n_c}{N}\,\mathrm{AUC}_c.
					\end{equation*}
					It is computed via \path{roc_auc_score(y_test, y_prob, multi_class='ovr', average='weighted')}, where \path{y_prob} is the probability matrix returned by \path{predict_proba}. A score of $0.5$ indicates random performance; $1.0$ indicates perfect discrimination. AUC-ROC is threshold-independent and provides a complete picture of model discrimination, complementing accuracy \cite{ref81}, \cite{ref82}.

				\subsubsection{Two-Percentage-Point Performance-Parity Threshold}
					Objective~4 operationalizes ``retaining predictive performance'' through a two-percentage-point tolerance on weighted F1-score and weighted OvR AUC-ROC. This threshold is not arbitrary. In the comparative interpretable-ML literature, differences on the order of 1--3 percentage points are routinely treated as substantively equivalent on tabular classification tasks. Nori et al.'s \cite{ref15} InterpretML benchmark reports EBM accuracy within 1--2~pp of XGBoost and random forests on several UCI tabular datasets and concludes that no meaningful performance loss is incurred by choosing the interpretable model. Caruana et al.\ \cite{ref50} report comparable parity (generalized additive model vs.\ boosted trees) on pneumonia-risk and hospital-readmission prediction. Rudin \cite{ref51} argues more generally that on structured tabular data the empirical performance gap between glass-box and black-box models is routinely small enough to be dominated by the interpretability benefit. A two-percentage-point threshold is therefore on the strict side of what this literature treats as practically equivalent, and passing it is a conservative bar for the claim that ``native interpretability is retained without a substantive loss.'' The threshold is also defensible at a defense-panel level: a gap of $\leq 2$~pp on a 1{,}000-sample held-out test set corresponds to a difference of at most 20 reclassifications, which is within the stability band of a single 80/20 split estimate \cite{ref68}.

		
		\section{RESULTS AND DISCUSSION}
		This chapter reports the results of the study and interprets them against the four objectives stated in the Introduction. The chapter is organized to mirror the sequence in which the objectives were executed, so that each section answers exactly one objective and presents the evidence for whether that objective was achieved. The first two subsections address the data-preparation objectives, documenting the synthetic dataset that was produced (Objective~1) and the rule-based labeling process that assigned products to each profile (Objective~2). The subsequent subsections address the model objectives, reporting the held-out test-set performance of both the Explainable Boosting Machine (EBM) and the XGBoost benchmark (Objective~3) and performing the formal two-percentage-point threshold analysis between them (Objective~4). All quantitative results were obtained from evaluation on the stratified 20\% held-out test set described in the Methodology, using scikit-learn \cite{ref74} for every metric computation. Values are quoted to four decimal places to match the precision of the saved model artifacts; percentages are rounded to two decimal places for readability.

			\subsection{Synthetic Dataset Generation}
				\subsubsection{Descriptive Statistics and Source Alignment}
					The descriptive statistics of the continuous features, computed directly from the generated dataset of $n = 5{,}000$ profiles, are reported in Table~\ref{tab:desc-stats}.

					\begin{table}[h]
					\vspace{4ex}
					\centering
					\caption{Descriptive Statistics of Continuous Features in the Synthetic Dataset ($n = 5{,}000$).}
					\label{tab:desc-stats}
					\scriptsize
					\begin{tabular}{|p{1.5cm}|r|r|r|r|r|r|}
						\hline
						\textbf{Feature} & \textbf{Min} & \textbf{Max} & \textbf{Mean} & \textbf{Median} & \textbf{Std Dev} & \textbf{Skewness} \\
						\hline
						age & 18 & 65 & 33.64 & 31 & 11.06 & 0.76 \\
						monthly\_\newline income & 252.48 & 175{,}672 & 22{,}922 & 18{,}035 & 20{,}874 & 1.91 \\
						monthly\_\newline expenses & 159.14 & 163{,}233 & 18{,}612 & 14{,}307 & 17{,}411 & 2.04 \\
						existing\_\newline savings & 0.00 & 796{,}153 & 25{,}850 & 11{,}099 & 43{,}607 & 4.66 \\
						num\_\newline dependents & 0 & 5 & 1.57 & 1 & 1.37 & 0.58 \\
						digital\_\newline savviness & 1 & 5 & 3.42 & 4 & 1.19 & $-0.42$ \\
						has\_bank\_\newline account & 0 & 1 & 0.488 & 0 & 0.500 & 0.05 \\
						has\_ewallet & 0 & 1 & 0.604 & 1 & 0.489 & $-0.43$ \\
						receives\_\newline remittance & 0 & 1 & 0.230 & 0 & 0.421 & 1.26 \\
						disposable\_\newline income & 18.94 & 57{,}714 & 4{,}310 & 2{,}975 & 4{,}597 & 2.56 \\
						expense\_ratio & 0.52 & 0.98 & 0.80 & 0.81 & 0.09 & $-0.32$ \\
						savings\_to\_\newline income\_ratio & 0.00 & 40.03 & 1.29 & 0.76 & 1.70 & 4.68 \\
						income\_per\_\newline dependent & 94.79 & 159{,}354 & 11{,}195 & 7{,}357 & 12{,}609 & 3.34 \\
						\hline
					\end{tabular}

					\vspace{0.5em}
					\raggedright\scriptsize \emph{Note.} Monetary values are in Philippine pesos (PHP). Skewness computed as the adjusted Fisher--Pearson coefficient. Income, expenses, and savings display the expected heavy right-skew of household-economics data \cite{ref54}.
					\vspace{4ex}
					\end{table}

					The categorical features in the generated dataset are distributed as follows. The \emph{employment\_status} feature is distributed as Employed (40.1\%), Self-employed (21.8\%), Unemployed (14.7\%), Freelancer (12.7\%), and Student (10.7\%), reflecting the working-age labor-participation patterns reported by PSA LFS \cite{ref41} together with an intentionally elevated Unemployed rate to ensure sufficient coverage of low-income profiles in the downstream model. The \emph{education} feature is distributed as High School (36.0\%), College (33.1\%), Elementary (12.9\%), Vocational (12.6\%), and Graduate (5.3\%), aligned with the PSA 2020 Census educational-attainment distribution \cite{ref39} and the PSA LFS attainment-by-age tabulations \cite{ref41}. The \emph{location\_type} feature is distributed as Urban (45.2\%), Rural (41.8\%), and Metro Manila (13.0\%), closely matching the 13\% Metro Manila share reported in the PSA 2020 Census \cite{ref39}. The \emph{primary\_ewallet} feature is distributed as Neither (39.6\%), GCash (32.0\%), Maya (16.1\%), and Both (12.3\%); the GCash--Maya leadership ratio of roughly 2:1 is consistent with provider-disclosed active-user shares and BSP FIS findings on e-money penetration \cite{ref9}, \cite{ref35}, and the Neither share aligns with the 39.6\% complement of the 60.4\% \emph{has\_ewallet} ownership rate. The \emph{savings\_goal} feature is distributed as Emergency Fund (23.4\%), Education (14.9\%), General Savings (14.8\%), Home/Property (13.3\%), Retirement (12.2\%), Business Capital (10.8\%), and Travel (10.6\%); the dominance of Emergency Fund is consistent with the BSP FIS finding that a majority of Filipino savers cite unforeseen expenses as their primary motivation \cite{ref35}. The \emph{risk\_tolerance} feature is distributed as Moderate (43.7\%), Conservative (36.1\%), and Aggressive (20.2\%), with the conservative-to-moderate majority consistent with the Philippine retail-investor profile documented in BSP consumer surveys. The \emph{investment\_horizon} feature is distributed as Medium-term (38.1\%), Short-term (32.3\%), and Long-term (29.6\%), reflecting the goal-conditional sampling described in the Methodology. The binary \emph{receives\_remittance} feature is realized at a rate of 23.0\%, in the same range as PSA FIES reporting on remittance-receiving households (approximately 22--25\% nationally) \cite{ref40} and the BSP FIS coverage of remittance channels \cite{ref35}.

					Table~\ref{tab:alignment} summarizes the quantitative alignment between key generated distributions and their Philippine source targets. The bank-account ownership rate of 48.80\% is below the BSP FIS 2021 headline of 56\% \cite{ref35} because the conditional-sampling rules apply income and employment penalties in \emph{generate\_behavioral\_features} alongside a deliberate oversampling of unemployed profiles, yielding a conservative estimate. The e-wallet ownership rate of 60.40\% is above the ``$>$50\%'' figure reported by BSP FIS 2021 \cite{ref35} and is consistent with the continued growth of GCash and Maya adoption between 2021 and 2024 \cite{ref42}. The remittance-recipient rate of 23.00\% is within the PSA FIES range of 22--25\% \cite{ref40}, and the College+Graduate share of 38.40\% is consistent with the PSA 2020 Census reporting that approximately one third of the working-age population has post-secondary attainment \cite{ref39}. Monthly income exhibits the expected lognormal right-skew (skewness $1.91$) \cite{ref54}, and existing savings are heavily right-skewed (skewness $4.66$, median PHP 11{,}099), matching the small savings buffers of the median Filipino household \cite{ref35}, \cite{ref55}.

					\begin{table}[h]
					\vspace{4ex}
					\centering
					\caption{Alignment of Generated Distributions with Philippine Source Targets.}
					\label{tab:alignment}
					\footnotesize
					\begin{tabular}{|p{4cm}|p{2cm}|p{2.5cm}|p{2cm}|}
						\hline
						\textbf{Variable} & \textbf{Source Target} & \textbf{Generated} & \textbf{Deviation} \\
						\hline
						Metro Manila share \cite{ref39} & 13.00\% & 13.00\% & $0.00$ pp \\
						Bank-account ownership \cite{ref35} & 56.00\% & 48.80\% & $-7.20$ pp \\
						E-wallet ownership \cite{ref35} & $>$50.00\% & 60.40\% & $>+10.40$ pp \\
						Remittance-receiving households \cite{ref40} & 22--25\% & 23.00\% & within range \\
						College + Graduate attainment \cite{ref39} & $\sim$33--40\% & 38.40\% & within range \\
						GCash:Maya primary share ratio \cite{ref9}, \cite{ref35} & $\sim$2:1 & 32.0\%:16.1\% (2.0:1) & consistent \\
						Income skewness (lognormal) \cite{ref54} & $>0$ (right-skew) & $1.91$ & consistent \\
						Savings skewness (heavy-tail) \cite{ref55} & $>0$ (right-skew) & $4.66$ & consistent \\
						Household size mean \cite{ref39} & $\sim$4.1 persons & $1.57$ dependents$^\dagger$ & see note \\
						\hline
					\end{tabular}

					\vspace{0.5em}
					\raggedright\scriptsize \emph{Note.} Deviations in pp indicate direction and magnitude of mismatch. $^\dagger$The \emph{num\_dependents} feature counts dependents of the individual respondent, not total household members; the PSA household-size figure of $\sim$4.1 corresponds to $\sim$3.1 dependents per household head, and the generated mean of 1.57 reflects the conditional distribution including unmarried young adults and students.
					\vspace{4ex}
					\end{table}

					All three success criteria were met: $n = 5{,}000$ profiles; 16 input features spanning 5 demographic (age, location, employment, education, num\_dependents), 3 financial (income, expenses, savings), 5 behavioral (digital savviness, has\_bank\_account, has\_ewallet, primary\_ewallet, receives\_remittance), and 3 preference (savings\_goal, risk\_tolerance, investment\_horizon) variables; and Philippine-grounded parameterization for every marginal and conditional distribution (Table~\ref{tab:features}). The approach follows Assefa et al.\ \cite{ref46} and Jordon et al.\ \cite{ref83}, who endorse parametric synthetic data under the privacy and data-access constraints imposed by the Philippine Data Privacy Act \cite{ref45} and the Financial Consumer Protection Act \cite{ref44}. The principal limitation, already acknowledged in the Methodology, is that the generator encodes selected conditional dependencies rather than a learned joint distribution of the kind provided by the Synthetic Data Vault \cite{ref29} or CTGAN \cite{ref34}.

			\subsection{Rule-Based Dataset Labeling}
				Eight product labels were assigned to every generated profile through the multi-factor additive scoring system specified in Section~\ref{sec:label-gen}, which accumulates bonuses across eleven scoring dimensions (platform attachment, digital preference, income level, existing savings, risk tolerance, investment horizon, savings goal, location, age, education, and remittance receipt) and selects the argmax product, with 8\% top-3 stochastic substitution to prevent deterministic labels. The ten rules (R1--R10) were detailed in the Methodology.

				\subsubsection{Class Distribution and Stratified Split}
					The resulting class distribution is reported in Table~\ref{tab:class-dist}.

					\begin{table}[h]
					\vspace{4ex}
					\centering
					\caption{Distribution of Product Labels in the Synthetic Dataset ($n = 5{,}000$).}
					\label{tab:class-dist}
					\footnotesize
					\begin{tabular}{|l|p{2.5cm}|l|r|r|}
						\hline
						\textbf{Product Key} & \textbf{Product Name} & \textbf{Provider} & \textbf{Count} & \textbf{\%} \\
						\hline
						gcash\_gsave & GSave (CIMB) & GCash & 1{,}737 & 34.74 \\
						bpi\_save\_up & BPI \#SaveUp & BPI & 805 & 16.10 \\
						maya\_personal\_goals & Maya Personal Goals & Maya & 639 & 12.78 \\
						gcash\_ginvest & GInvest (UITFs) & GCash & 562 & 11.24 \\
						maya\_savings & Maya Savings & Maya & 424 & 8.48 \\
						bpi\_time\_deposit & BPI Plan Ahead Time Deposit & BPI & 392 & 7.84 \\
						bpi\_savings & BPI Regular Savings & BPI & 320 & 6.40 \\
						maya\_time\_deposit & Maya Time Deposit Plus & Maya & 121 & 2.42 \\
						\hline
						\textbf{Total} & & & \textbf{5{,}000} & \textbf{100.00} \\
						\hline
					\end{tabular}
					\vspace{4ex}
					\end{table}

					The distribution is imbalanced: the largest class (\emph{gcash\_gsave} at 34.74\%) is roughly 14.4$\times$ more frequent than the smallest (\emph{maya\_time\_deposit} at 2.42\%). This imbalance widened relative to the previous-iteration generator after \emph{primary\_ewallet} was introduced as an observable platform-attachment signal: \emph{gcash\_gsave} consolidates both the population's largest e-wallet primary-user segment (GCash, 32.0\%) and a portion of the Both segment (12.3\%), and is also the only no-minimum, no-eligibility-friction option for low-income profiles. The pattern is emergent from the generator's demographic mixture (predominantly young, digitally engaged, low-to-moderate-income profiles with Emergency Fund as the dominant goal) interacting with the product-scoring rules; it mirrors the real Philippine market structure in which mass-market digital wallets dominate over traditional deposit products \cite{ref35}, \cite{ref42}. Rather than resampling via SMOTE \cite{ref47}, which would distort the population mixture, the imbalance is handled at evaluation time by weighted averaging \cite{ref58}, \cite{ref78}. The verified per-class distribution after the stratified 80/20 split \cite{ref59} is reported in Table~\ref{tab:split}; both the EBM and XGBoost pipelines consume identical row indices (same \emph{random\_state}$=42$, same stratification vector), differing only in feature preprocessing.

					\begin{table}[h]
					\vspace{4ex}
					\centering
					\caption{Class Distribution in Training and Test Sets After Stratified 80/20 Split.}
					\label{tab:split}
					\footnotesize
					\begin{tabular}{|p{2.5cm}|r|r|r|r|c|}
						\hline
						\textbf{Product} & \textbf{Train N} & \textbf{Train \%} & \textbf{Test N} & \textbf{Test \%} & \textbf{Diff} \\
						\hline
						gcash\_gsave & 1{,}389 & 34.73 & 348 & 34.80 & 0.07 pp \\
						bpi\_save\_up & 644 & 16.10 & 161 & 16.10 & 0.00 pp \\
						maya\_personal\_\newline goals & 511 & 12.78 & 128 & 12.80 & 0.02 pp \\
						gcash\_ginvest & 450 & 11.25 & 112 & 11.20 & 0.05 pp \\
						maya\_savings & 339 & 8.48 & 85 & 8.50 & 0.02 pp \\
						bpi\_time\_deposit & 314 & 7.85 & 78 & 7.80 & 0.05 pp \\
						bpi\_savings & 256 & 6.40 & 64 & 6.40 & 0.00 pp \\
						maya\_time\_deposit & 97 & 2.43 & 24 & 2.40 & 0.03 pp \\
						\hline
						\textbf{Total} & \textbf{4{,}000} & 100.00 & \textbf{1{,}000} & 100.00 & \\
						\hline
					\end{tabular}

					\vspace{0.5em}
					\raggedright\scriptsize \emph{Note.} Stratified split parameters: \path{test_size=0.2}, \path{random_state=42}, \path{stratify=y} \cite{ref59}, \cite{ref74}.
					\vspace{4ex}
					\end{table}

					The maximum class-percentage deviation between the train and test sets is less than $0.10$ pp across all eight classes, confirming that stratification preserved the full-dataset class priors.

				\subsubsection{Discussion of Objective 2 Achievement and Circular-Reasoning Risk}
					The eight rules were traced directly to the official product documentation of GCash \cite{ref6}, Maya, and BPI \cite{ref37}, and rule-based ground truth is established practice in financial classification when real outcome labels are unavailable \cite{ref46}, \cite{ref56}, \cite{ref57}.

					A fundamental epistemological concern arises from the study design: because labels are generated by rules and the classifier is evaluated against those same labels, a skeptical reviewer may ask whether the model is genuinely learning financial-suitability patterns or merely memorizing the labeling rules. Three structural features of the generator attenuate this risk. First, the additive scoring combines eleven weakly aligned dimensions (platform attachment, digital, income, savings, risk, horizon, goal, location, age, education, remittance) whose joint effect is not a clean decision surface---multiple rules overlap on many profiles (e.g., a digitally savvy low-savings Emergency-Fund saver with \emph{primary\_ewallet} $=$ Maya satisfies both R1 and R3), so the argmax decision requires the model to learn relative magnitudes, not a single threshold. Second, an intentional 8\% top-3 label-noise substitution injects stochasticity into the training signal, so even a model that perfectly recovered the scoring arithmetic would still face an irreducible labeling error of at least eight percentage points; this places a hard upper bound on attainable accuracy below 1.0 regardless of model capacity. Third, if the EBM were simply ``reading off'' the rules it should achieve near-100\% F1 rather than the 0.8711 weighted F1 reported in the following subsection, with persistent residuals on the smallest class (\emph{maya\_time\_deposit} F1 $=0.5581$); the observed gap and its concentration in low-support classes is consistent with genuine generalization under noise rather than memorization. These mitigations are partial, not definitive. The canonical test---training on this synthetic dataset and evaluating on real Filipino banked-customer application outcomes---remains future work and is the condition under which external validity can be established \cite{ref46}.

					Real-world labels would additionally reflect credit history, income-verification documentation, Know-Your-Customer (KYC) checks, and behavioral signals not modeled here \cite{ref46}, which is why the dataset is positioned as a research prototype and not a deployment artifact.

			\subsection{Model Training and Performance Evaluation}
				Objective~3 was achieved. Both the EBM and the XGBoost benchmark were trained on the 4{,}000-row training set and evaluated on the 1{,}000-row stratified held-out test set described above, using the five metrics and the scikit-learn function calls enumerated in the Methodology \cite{ref74}.

				\subsubsection{Training Configuration}
					The EBM and XGBoost hyperparameters were read verbatim from \emph{notebooks/train\_ebm.py} and \emph{notebooks/train\_model.py} and are summarized in Tables~\ref{tab:ebm-hp} and~\ref{tab:xgb-hp}. The XGBoost values are the author-selected defaults documented in the Methodology rather than the output of a systematic grid search; a formal hyperparameter-tuning comparison is deferred to future work (Recommendation R4) and the results below should therefore be read as a comparison of a thoughtfully-configured EBM against an untuned XGBoost baseline. Preprocessing pipelines remain as specified in the Methodology: the EBM consumes the 20-column DataFrame natively (\emph{category} dtype for the seven categoricals \emph{employment\_status}, \emph{education}, \emph{location\_type}, \emph{primary\_ewallet}, \emph{savings\_goal}, \emph{risk\_tolerance}, and \emph{investment\_horizon}), while XGBoost applies \emph{OrdinalEncoder} and \emph{StandardScaler} \cite{ref74}. No class weights were used; imbalance is handled at evaluation time via weighted averaging.

					\begin{table}[h]
					\vspace{4ex}
					\centering
					\caption{EBM Hyperparameter Configuration (verbatim from \emph{train\_ebm.py}).}
					\label{tab:ebm-hp}
					\footnotesize
					\begin{tabular}{|l|l|p{4cm}|}
						\hline
						\textbf{Hyperparameter} & \textbf{Value} & \textbf{Description} \\
						\hline
						\emph{max\_bins} & 256 & Maximum bin count per feature for histogramming \\
						\emph{interactions} & 10 & Number of pairwise-interaction terms (GA$^2$M) \\
						\emph{outer\_bags} & 8 & Bagging rounds for variance reduction \\
						\emph{inner\_bags} & 0 & Disabled; bagging only at outer level \\
						\emph{learning\_rate} & 0.01 & Boosting step size \\
						\emph{max\_rounds} & 5000 & Boosting round cap before early stopping \\
						\emph{min\_samples\_leaf} & 2 & Minimum samples per leaf in each booster \\
						\emph{random\_state} & 42 & Reproducibility seed \\
						\emph{feature\_names} & 20 entries & Full ordered feature list incl.\ derived \\
						Others & defaults & InterpretML built-in defaults \cite{ref15} \\
						\hline
					\end{tabular}
					\vspace{4ex}
					\end{table}

					\begin{table}[h]
					\vspace{4ex}
					\centering
					\caption{XGBoost Hyperparameter Configuration (verbatim from \emph{train\_model.py}).}
					\label{tab:xgb-hp}
					\footnotesize
					\begin{tabular}{|l|p{2.2cm}|p{2.7cm}|}
						\hline
						\textbf{Hyperparameter} & \textbf{Value} & \textbf{Description} \\
						\hline
						\emph{n\_estimators} & 300 & Number of boosting trees \\
						\emph{max\_depth} & 6 & Maximum tree depth \\
						\emph{learning\_rate} & 0.1 & Boosting step size \\
						\emph{subsample} & 0.8 & Row subsample per tree \\
						\emph{colsample\_bytree} & 0.8 & Column subsample per tree \\
						\emph{min\_child\_weight} & 3 & Minimum child sum of instance weight \\
						\emph{gamma} & 0.1 & Minimum loss reduction for a split \\
						\emph{reg\_alpha} & 0.1 & L$_1$ regularization on weights \\
						\emph{reg\_lambda} & 1.0 & L$_2$ regularization on weights \\
						\emph{objective} & \emph{multi:softprob} & Multiclass softmax probabilistic objective \\
						\emph{num\_class} & 8 & Number of target classes \\
						\emph{eval\_metric} & \emph{mlogloss} & Multi-class log loss monitor \\
						\emph{random\_state} & 42 & Reproducibility seed \\
						\hline
					\end{tabular}
					\vspace{4ex}
					\end{table}

					Both models were trained on the identical 4{,}000-row training set produced by the same \emph{train\_test\_split(..., test\_size=0.2, random\_state=42, stratify=y)} call, so that any performance difference is attributable to the choice of model rather than to the training sample. A single train/test split was used for the five-metric headline evaluation due to compute constraints; this is a known limitation and Dem\v{s}ar \cite{ref68} recommends multi-seed or multi-dataset replication before strong superiority claims are made. Future work (Recommendation R4) will report confidence intervals on the EBM--XGBoost gap across multiple random seeds.

					As a stability check on the EBM side, stratified five-fold cross-validation was run on the full 5{,}000-row dataset using the same \emph{random\_state=42} seed (see \emph{notebooks/train\_ebm.py} and \emph{models/model\_metadata.json}). The five per-fold accuracies were $\{0.868, 0.868, 0.860, 0.877, 0.877\}$, giving a mean of $0.870$ and a sample standard deviation of approximately $0.0064$; the held-out test-set accuracy of $0.8730$ therefore sits within roughly $0.5\sigma$ of the cross-validation mean, well inside the cross-fold stability band, indicating the headline test-set result is not an artifact of a single fortuitous split. A matching cross-validation run for the XGBoost benchmark was not performed and is explicitly scoped to Recommendation~R4, so the number reported here is an EBM-side stability indicator rather than a side-by-side cross-validation comparison between the two models.

				\subsubsection{Performance Results}
					The five-metric performance of both models on the stratified held-out test set is reported in Table~\ref{tab:perf}. The difference column reports EBM minus XGBoost in percentage points, with a positive value indicating that EBM is superior on that metric.

					\begin{table}[h]
					\vspace{4ex}
					\centering
					\caption{Model Performance on the Stratified Held-Out Test Set ($n_{\mathrm{test}} = 1{,}000$, 8 Classes). Both learned models are compared against a majority-class (Maj.) baseline that always predicts \emph{gcash\_gsave} (the dominant class at 34.8\% support).}
					\label{tab:perf}
					\footnotesize
					\begin{tabular}{|p{2.5cm}|c|c|c|p{1.3cm}|}
						\hline
						\textbf{Metric} & \textbf{EBM} & \textbf{XGBoost} & \textbf{Maj.} & \textbf{EBM $-$ XGB (pp)} \\
						\hline
						Accuracy & $0.8730$ & $0.8150$ & $0.3480$ & $+5.80$ \\
						Weighted Precision & $0.8720$ & $0.8138$ & $0.1211$ & $+5.82$ \\
						Weighted Recall & $0.8730$ & $0.8150$ & $0.3480$ & $+5.80$ \\
						Weighted F1-Score & $0.8711$ & $0.8118$ & $0.1797$ & $+5.93$ \\
						Weighted OvR AUC-ROC & $0.9761$ & $0.9671$ & $0.5000$ & $+0.90$ \\
						\hline
					\end{tabular}

					\vspace{0.5em}
					\raggedright\scriptsize \emph{Note.} EBM $-$ XGB (pp): EBM minus XGBoost, in percentage points. The majority-class column (Maj.) shows the metric values that would be obtained by a trivial classifier that always predicts \emph{gcash\_gsave}; by construction its accuracy equals the prior of the majority class, its weighted OvR AUC-ROC is 0.5 (no discrimination), and both learned models exceed it by approximately 47--52\,pp on accuracy and 63--69\,pp on weighted F1. EBM outperforms XGBoost on every reported metric under identical train/test splits and evaluation protocol. Metrics computed via scikit-learn \cite{ref74}.
					\vspace{4ex}
					\end{table}

					The per-class classification reports for the two models on the same test set are reported in Tables~\ref{tab:ebm-per-class} and~\ref{tab:xgb-per-class}.

					\begin{table}[h]
					\vspace{4ex}
					\centering
					\caption{Per-Class Classification Report for the EBM Model on the Held-Out Test Set.}
					\label{tab:ebm-per-class}
					\footnotesize
					\begin{tabular}{|l|c|c|c|c|}
						\hline
						\textbf{Product} & \textbf{Precision} & \textbf{Recall} & \textbf{F1} & \textbf{Support} \\
						\hline
						bpi\_save\_up & 0.8516 & 0.8199 & 0.8354 & 161 \\
						bpi\_savings & 0.8704 & 0.7344 & 0.7966 & 64 \\
						bpi\_time\_deposit & 0.8235 & 0.8974 & 0.8589 & 78 \\
						gcash\_ginvest & 0.9238 & 0.8661 & 0.8940 & 112 \\
						gcash\_gsave & 0.8859 & 0.9598 & 0.9214 & 348 \\
						maya\_personal\_goals & 0.8689 & 0.8281 & 0.8480 & 128 \\
						maya\_savings & 0.9036 & 0.8824 & 0.8929 & 85 \\
						maya\_time\_deposit & 0.6316 & 0.5000 & 0.5581 & 24 \\
						\hline
						macro avg & 0.8449 & 0.8110 & 0.8257 & 1{,}000 \\
						weighted avg & 0.8720 & 0.8730 & 0.8711 & 1{,}000 \\
						\hline
					\end{tabular}
					\vspace{4ex}
					\end{table}

					\begin{table}[h]
					\vspace{4ex}
					\centering
					\caption{Per-Class Classification Report for the XGBoost Benchmark on the Held-Out Test Set.}
					\label{tab:xgb-per-class}
					\footnotesize
					\begin{tabular}{|l|c|c|c|c|}
						\hline
						\textbf{Product} & \textbf{Precision} & \textbf{Recall} & \textbf{F1} & \textbf{Support} \\
						\hline
						bpi\_save\_up & 0.7799 & 0.7702 & 0.7750 & 161 \\
						bpi\_savings & 0.8333 & 0.7031 & 0.7627 & 64 \\
						bpi\_time\_deposit & 0.8400 & 0.8077 & 0.8235 & 78 \\
						gcash\_ginvest & 0.8587 & 0.7054 & 0.7745 & 112 \\
						gcash\_gsave & 0.8389 & 0.9425 & 0.8877 & 348 \\
						maya\_personal\_goals & 0.7287 & 0.7344 & 0.7315 & 128 \\
						maya\_savings & 0.8780 & 0.8471 & 0.8623 & 85 \\
						maya\_time\_deposit & 0.5556 & 0.4167 & 0.4762 & 24 \\
						\hline
						macro avg & 0.7891 & 0.7409 & 0.7617 & 1{,}000 \\
						weighted avg & 0.8138 & 0.8150 & 0.8118 & 1{,}000 \\
						\hline
					\end{tabular}
					\vspace{4ex}
					\end{table}

					Three qualitative patterns emerge from the per-class reports and the confusion matrices. First, the dominant class \emph{gcash\_gsave} and the high-signal investment class \emph{gcash\_ginvest} are the best-discriminated for the EBM (F1 $=0.9214$ and $0.8940$ respectively); this is consistent with the aggressive-risk, long-horizon scoring pattern in Rule~R2 plus the strong \emph{primary\_ewallet} $=$ GCash platform-attachment signal, which together create a distinctive feature-space signature. Second, the previously pathological \emph{maya\_savings} class is now well-recovered by both models (EBM F1 $=0.8929$; XGBoost F1 $=0.8623$), confirming that the introduction of \emph{primary\_ewallet} successfully resolved the structural overlap with \emph{gcash\_gsave} flagged by the defense panel: the EBM's confusion matrix shows only 6 of the 85 true \emph{maya\_savings} examples now misclassified as \emph{gcash\_gsave}, down from 117 of 133 in the prior generator iteration. Third, the residual weakness is now concentrated on the smallest class \emph{maya\_time\_deposit} (EBM F1 $=0.5581$, support $=24$), which is the expected pattern under class imbalance and is consistent with low-support classes in the imbalanced-multiclass literature \cite{ref58}; \emph{bpi\_savings} also drops to recall $0.7344$ on its 64-sample test partition, reflecting the small-class--limited evidence available to the model.

					The confusion matrix for the EBM is reported in Fig.~\ref{fig:cm-ebm} (diagonal cells for perfect prediction, off-diagonal cells for misclassification). After the introduction of \emph{primary\_ewallet}, no single off-diagonal pair dominates as the previous \emph{maya\_savings}$\to$\emph{gcash\_gsave} cluster did. The largest remaining residual confusions cluster around the dominant \emph{gcash\_gsave} class: \emph{bpi\_save\_up}$\to$\emph{gcash\_gsave} (11), \emph{maya\_personal\_goals}$\to$\emph{gcash\_gsave} (11), and within-BPI \emph{bpi\_savings}$\to$\emph{bpi\_save\_up} (11)---all reflecting the residual draw of the highest-prior class on borderline profiles. The confusion matrix for the XGBoost benchmark (Fig.~\ref{fig:cm-xgb}) shows the same overall pattern but spreads additional errors more evenly across smaller classes, consistent with XGBoost's lower per-class F1 on \emph{maya\_personal\_goals} ($0.7315$ vs.\ EBM $0.8480$) and \emph{bpi\_save\_up} ($0.7750$ vs.\ EBM $0.8354$).

					\begin{figure}[ht]
					\begin{center}
						\vspace{4ex}
						\includegraphics[width=1\textwidth]{cm_ebm.png}
						\label{fig:cm-ebm}
						\caption{Confusion matrix for the EBM model on the stratified held-out test set. Rows represent true product labels; columns represent predicted labels. Darker cells indicate higher prediction frequency.}
					\end{center}
					\end{figure}

					\begin{figure}[ht]
					\begin{center}
						\vspace{4ex}
						\includegraphics[width=1\textwidth]{cm_xgb.png}
						\label{fig:cm-xgb}
						\caption{Confusion matrix for the XGBoost benchmark model on the same stratified held-out test set. Rows represent true product labels; columns represent predicted labels.}
					\end{center}
					\end{figure}

				\subsubsection{Discussion of Performance Metrics}
					Both models substantially exceed the uniform-random baseline of $1/8 = 0.125$ and the majority-class baseline of $0.348$ (always predict \emph{gcash\_gsave}), delivering roughly 2.5$\times$ the majority-class accuracy. The EBM outperforms XGBoost on every metric by $+0.90$ to $+5.93$ pp (Table~\ref{tab:perf}); the same-sign advantage on both precision and recall indicates that the EBM's gain is not a favorable false-positive/false-negative trade-off but a genuine improvement in classification quality. Weighted F1 ($0.8711$ vs.\ $0.8118$) is the headline metric because it is robust to the class imbalance documented above \cite{ref58}, \cite{ref78}.

					Contextualizing the EBM's weighted F1 of $0.8711$ requires acknowledging that no prior study reports multi-class Philippine fintech product-recommendation metrics on a comparable 8-class task. For reference, West \cite{ref56} reports F1 in the 0.70--0.85 range on \emph{binary} credit-scoring with real labels; the 8-class setting with synthetic labels and 8\% label noise is a strictly harder problem, so the present result sits at the upper end of that benchmark range despite the harder problem class. Bussmann et al.\ \cite{ref14} report accuracies in the 0.60--0.75 range for explainable credit-risk classification on real European data, placing the present synthetic result above that band but with the caveat that real-data validation remains future work (Recommendation~R1). No comparable Philippine-fintech benchmark exists, which is itself a gap this study begins to fill.

					The weighted OvR AUC-ROC of $0.9761$ (EBM) and $0.9671$ (XGBoost) indicates that both models rank the correct product very high in the predicted probability distribution, with the EBM in particular reaching the high-discrimination band of Fawcett's \cite{ref81} informal AUC-ROC heuristic ($>0.9$ excellent; $0.8$--$0.9$ good). The closer agreement between accuracy ($0.873$) and AUC-ROC ($0.976$) on the new feature set, relative to the previous-iteration gap, reflects the disambiguation effect of \emph{primary\_ewallet}: when the dominant \emph{maya\_savings}-vs-\emph{gcash\_gsave} confusion is resolved, the model's argmax prediction agrees with its probability ranking far more often.

			\subsection{EBM vs.\ XGBoost: Interpretability-Performance Trade-off Analysis}
				The formal two-percentage-point threshold analysis is reported in Table~\ref{tab:threshold}: the EBM exceeded XGBoost on weighted F1 by $+5.93$ pp and on weighted OvR AUC-ROC by $+0.90$ pp, so the stronger ``greater than'' clause of Objective~4 is satisfied on both headline metrics. The weaker ``within 2~pp'' clause does not need to be invoked, because on this dataset the EBM is on the favorable side of the band for both headline metrics. The F1 advantage is now substantively meaningful (well beyond the two-percentage-point tolerance), driven primarily by the EBM's better recovery of the previously confused \emph{maya\_savings} class and its tighter handling of the GCash--Maya digital-savings boundary; the AUC-ROC gap is small ($+0.90$ pp), within the band, and reflects that both models now achieve near-saturation discrimination on this task. Because AUC-ROC is threshold-independent \cite{ref81}, \cite{ref82}, the EBM's advantage reflects superior ranking quality across the full probability range rather than a favorable operating point.

				\begin{table}[h]
				\vspace{4ex}
				\centering
				\caption{Two-Percentage-Point Threshold Analysis: EBM vs.\ XGBoost on the Held-Out Test Set.}
				\label{tab:threshold}
				\footnotesize
				\begin{tabular}{|p{1.7cm}|c|c|c|c|}
					\hline
					\textbf{Metric} & \textbf{EBM} & \textbf{XGBoost} & \textbf{Diff (pp)} & \textbf{Criterion Met?} \\
					\hline
					Weighted F1 & $0.8711$ & $0.8118$ & $+5.93$ & EBM superior \\
					Weighted AUC-ROC & $0.9761$ & $0.9671$ & $+0.90$ & EBM superior \\
					\hline
				\end{tabular}

				\vspace{0.5em}
				\raggedright\scriptsize \emph{Note.} Objective~4 states that EBM's weighted F1 and AUC-ROC must be within 2~pp of, or greater than, XGBoost. ``EBM superior'' indicates the stronger ``greater than'' clause is satisfied. The $\leq 2$pp column is marked only when the EBM is behind by $\leq 2$pp (not applicable when EBM is ahead).
				\vspace{4ex}
				\end{table}

			\subsection{Expert Review by a Practising Financial Advisor}
				Written feedback was solicited from Mr.\ Michael Chan, a Registered Financial Planner in the Philippines. The review focused specifically on the product-recommendation logic documented in the Methodology (Rules R1--R10) and on the per-class results reported above.

				The reviewer affirmed three substantive points. First, the derivation of the eight product-eligibility rules from each platform's public documentation was assessed as a reasonable approach for a recommendation prototype, given that the rules reflect each product's stated features (minimum deposit, withdrawal conditions, risk level, and intended user profile); the suitability dimensions used in the study's additive scoring system (liquidity, risk tolerance, investment horizon, and minimum balance) were identified as the same dimensions used in standard Philippine financial-advisory practice. Second, the inter-class confusion between \emph{gcash\_gsave} and \emph{maya\_savings} reported in the prior generator iteration was characterised as a faithful reflection of real-market conditions rather than a modeling defect, because the two products share very similar structural attributes (low entry barriers, competitive interest rates, high liquidity, minimal risk) and are differentiated in practice by ecosystem preference, user-interface familiarity, and promotional campaigns; this practitioner framing retroactively supports the panel-driven introduction of \emph{primary\_ewallet} as the appropriate observable discriminator. Third, the methodological combination of machine learning with transparent rule-based logic was endorsed as appropriate for financial applications in which explainability and user trust are material considerations.

				The reviewer raised one substantive gap: although the study captures stated suitability inputs (savings goal, risk tolerance, investment horizon, and a snapshot financial profile), it does not capture \emph{behavioural} or \emph{ongoing-suitability} signals---recurring-deposit patterns, withdrawal frequency, balance volatility, and similar longitudinal transactional behaviours that a real advisory engagement would normally incorporate. This gap is acknowledged explicitly in the limitations subsection and is addressed as a directed extension in the next chapter. It cannot be closed through synthetic data alone, because it presupposes access to transactional records under appropriate data-sharing agreements with BSP-supervised financial institutions; it is therefore best treated as a downstream extension of the real-world data validation recommendation.

		\section{SUMMARY AND CONCLUSION}
		This chapter synthesizes the evidence presented in the Results and Discussion chapter into a coherent closing account of the study. It consolidates the empirical findings, states the contributions that the study makes to the Philippine explainable-fintech literature, and sets out recommendations for future research and for practical deployment.

			\subsection{Summary of Findings}
				\subsubsection{Objectives Achievement}
					The status of each of the four objectives is summarized in Table~\ref{tab:obj}.

					\begin{table}[ht]
					\vspace{4ex}
					\centering
					\caption{Summary of Objectives Achievement.}
					\label{tab:obj}
					\footnotesize
					\begin{tabular}{|c|p{4.5cm}|p{6.5cm}|c|}
						\hline
						\textbf{\#} & \textbf{Objective} & \textbf{Key Result} & \textbf{Status} \\
						\hline
						1 & Generate synthetic dataset of $\geq 5{,}000$ profiles across 16 features, distributions from PSA/BSP sources & $n = 5{,}000$; 16 input features $+ 4$ derived; distributions parameterized from PSA FIES 2023, PSA LFS 2023, PSA 2020 Census, and BSP FIS 2021 & Achieved \\
						\hline
						2 & Label dataset via rule-based logic aligned to 8 product eligibility criteria & 8 product labels assigned via multi-factor additive scoring across 11 scoring dimensions; rules derived from GCash, Maya, and BPI official documentation; class distribution imbalanced (34.74\%--2.42\%) and handled via stratified split and weighted metrics & Achieved \\
						\hline
						3 & Train EBM and XGBoost, evaluate on 5 metrics on 20\% stratified test set & EBM: Accuracy $0.8730$, W-Prec $0.8720$, W-Rec $0.8730$, W-F1 $0.8711$, W-AUC $0.9761$; XGBoost: Accuracy $0.8150$, W-Prec $0.8138$, W-Rec $0.8150$, W-F1 $0.8118$, W-AUC $0.9671$. Per-class EBM F1 ranges from $0.5581$ (\emph{maya\_time\_deposit}, smallest class, support $24$) to $0.9214$ (\emph{gcash\_gsave}); previous \emph{maya\_savings} pathology resolved (F1 now $0.8929$). & Achieved \\
						\hline
						4 & EBM F1 and AUC within 2pp of or greater than XGBoost & $\Delta$W-F1 $= +5.93$ pp (EBM superior); $\Delta$W-AUC $= +0.90$ pp (EBM superior); criterion exceeded on both headline metrics & Achieved \\
						\hline
					\end{tabular}
					\vspace{4ex}
					\end{table}

				\subsubsection{Synthesized Findings}
					Taken together, the four objectives demonstrate the feasibility of an end-to-end transparent AI-driven savings-product recommender for Filipino fintech users: a Philippines-grounded synthetic dataset \cite{ref35}, \cite{ref39}, \cite{ref40}, \cite{ref41} assembled under the access constraints of the Data Privacy Act \cite{ref45}; a rule-based labeling system whose eligibility criteria trace to the official documentation of GCash \cite{ref6}, Maya, and BPI \cite{ref37}; and an EBM that exceeds a strong XGBoost benchmark on both headline metrics while preserving the native interpretability required by the Philippine Financial Consumer Protection Act \cite{ref44}. Because \emph{explain\_local} reflects the model's decision function exactly, the trustworthiness of the displayed \emph{key\_factors}, plain-language reasons, and counterfactuals rests on the model's own contributions rather than on post-hoc approximations \cite{ref15}, \cite{ref48}. These results address the problem framing of the Introduction: transparent, regulator-ready explanations are both feasible and empirically competitive with opaque alternatives.

			\subsection{Contributions of the Study}
				The study makes four contributions to the Philippine fintech and explainable-AI literature.

				\textit{Contribution~1---A Philippine-contextualized synthetic fintech dataset.} The study produced a reproducible synthetic dataset of 5{,}000 Filipino fintech user profiles parameterized jointly from PSA FIES 2023 \cite{ref40}, PSA LFS 2023 \cite{ref41}, PSA 2020 Census \cite{ref39}, and BSP Financial Inclusion Survey 2021 \cite{ref35}. To the best of the author's knowledge no prior publicly available synthetic dataset specifically models the financial-profile distribution of the Philippine retail savings market across these four sources simultaneously. Prior synthetic-data work in finance \cite{ref46}, \cite{ref83} and in tabular ML \cite{ref29}, \cite{ref34} has addressed either generic financial use cases or methodology rather than the Philippine-specific parameterization reported here. The generation script (with a fixed random seed \cite{ref60}), the 20-column preprocessed dataset, and the feature-to-source mapping in Table~\ref{tab:features} constitute a reusable research artifact for future Philippine fintech ML work.

				\textit{Contribution~2---Empirical evidence for EBM suitability on a Philippine fintech synthetic benchmark.} The study provides single-split empirical evidence, corroborated by a five-fold cross-validation stability check ($0.870 \pm 0.006$ accuracy), under a fixed random seed and without systematic XGBoost tuning, that the EBM attained a weighted F1 of $0.8711$ and a weighted OvR AUC-ROC of $0.9761$ on an 8-class Philippine savings-and-investment recommendation task, exceeding the XGBoost benchmark by $+5.93$ and $+0.90$ percentage points respectively (Table~\ref{tab:threshold}). This is a bounded contribution: the result is not yet validated on real Filipino banked-customer application data and the comparison would benefit from multi-seed and tuned-XGBoost replication \cite{ref68} before its magnitude is treated as established. Nevertheless, the direction of the effect is consistent with the growing body of intelligible-models evidence \cite{ref15}, \cite{ref48}, \cite{ref49}, \cite{ref50}, and supports the position that glass-box models based on shape functions are a defensible primary choice for Philippine fintech product-recommendation tasks where the Financial Products and Services Consumer Protection Act \cite{ref44} mandates regulatory transparency \cite{ref51}. This methodological choice has additionally received external practitioner endorsement: a practising Philippine financial planner characterised the combination of machine learning with transparent rule-based logic as appropriate for financial applications in which explainability and user trust are material considerations, complementing the quantitative argument with practitioner face-validity.

				\textit{Contribution~3---A working XAI prototype for Philippine fintech.} The study produced a functional web-based prototype that layers multiple complementary explanation modalities over a single EBM prediction: quantitative per-feature contributions from the EBM's native \emph{explain\_local} output \cite{ref15}; qualitative plain-language supporting and opposing reasons rendered from a deterministic template library; and actionable counterfactual ``what-if'' suggestions over the actionable-feature subset \cite{ref20}, \cite{ref86}. SHAP \cite{ref18} and LIME \cite{ref19} are retained in the offline benchmark pipeline for comparison against the XGBoost benchmark, producing a taxonomically complete set of the XAI modalities identified by Arrieta et al.\ \cite{ref17}. The contribution is an \emph{integration pattern}---a reusable architecture showing how these four explanation modalities can be combined behind a single fintech recommendation endpoint---rather than a validated claim that the pattern is comprehensible to end users; usability testing with Filipino participants is deferred to Recommendation R3. As such the contribution is to explainable-fintech \emph{system engineering} rather than to human--computer interaction.

				\textit{Contribution~4---A documented framework for Philippine fintech provider and product selection.} The study produced a systematic justification for selecting GCash, Maya Bank, and BPI as the three providers and eight specific products covered by the prototype, grounded in provider-level institutional evidence (transaction volume, user base, BSP licensing, and cross-provider coverage of the digital-mass-market / digital-bank / traditional-bank archetypes) and product-level suitability evidence (minimum balances, interest rates, liquidity terms, and BSP/PDIC-insurance status). This framework, summarized in the Methodology and supported by a standalone documentation artifact, is offered as a \emph{methodological} rather than a theoretical contribution: it gives future Philippine fintech studies a reusable, auditable structure for justifying a bounded product catalogue in the absence of a regulator-maintained canonical list, and it documents the specific BSP/PDIC/Universal-Banking institutional evidence that supports the three-archetype framing.

			\subsection{Recommendations}
				\subsubsection{Recommendations for Future Research}
					Of the five research recommendations listed below, Recommendation~R1 (real-world data validation) is the highest-priority next step: every other recommendation assumes, but cannot guarantee, that the synthetic-data results reported in this study transfer to real Filipino banked-customer profiles. R1 is therefore the prerequisite for any stronger claim about the FinXplain approach, and the remaining recommendations are best undertaken in parallel with, or subsequent to, an R1 validation.

					\textit{R1---Real-world data validation.} Future work should validate the FinXplain model on real Philippine bank-customer data obtained through ethical data-sharing agreements with BSP-supervised financial institutions or through a BSP-sponsored research program. The synthetic dataset, while grounded in national statistics, cannot fully replicate the joint-distribution feature correlations present in real customer profiles \cite{ref29}, \cite{ref34}, and training-on-synthetic / testing-on-real evaluation is the canonical methodology in the synthetic-data-in-finance literature for establishing external validity \cite{ref46}, \cite{ref83}. 

					\textit{R2---Expanded provider and product coverage.} The present study covers eight products from three providers. Future work should extend the catalogue to include additional BSP-licensed digital banks \cite{ref36}, \cite{ref38} such as Tonik Digital Bank, UNObank, and UnionDigital, as well as traditional universal banks such as BDO Unibank and UnionBank, and government-backed savings instruments such as the Pag-IBIG Modified Pag-IBIG 2 (MP2) program and Social Security System PESO Fund. Broader coverage would improve recommendation relevance across additional segments of the Philippine retail savings market and reduce the risk that the system systematically under-recommends products outside the current three-provider set.

					\textit{R3---Formal usability study with Filipino users.} A formal usability evaluation with Filipino participants---particularly older adults and users with limited financial literacy \cite{ref43}, \cite{ref76}---should be conducted. An established heuristic evaluation framework \cite{ref28} combined with task-based think-aloud sessions would test whether the EBM key-factor chart, the plain-language reasons, and the counterfactual suggestions are actually comprehensible to the target user population.

					\textit{R4---Tuned-XGBoost benchmark and multi-seed replication.} The XGBoost benchmark in this study uses author-selected default hyperparameters rather than the output of a systematic grid or Bayesian search. Future work should re-run the EBM-vs.-XGBoost comparison with a tuned XGBoost (e.g., via 5-fold cross-validated grid search over \emph{n\_estimators}, \emph{max\_depth}, \emph{learning\_rate}, \emph{min\_child\_weight}, \emph{gamma}, and the regularization terms) and across multiple random seeds, reporting confidence intervals on the EBM--XGBoost gap rather than a single-split point estimate \cite{ref68}. A matching cross-validation run on the XGBoost side, paired with the EBM stability check already reported above, would provide a side-by-side stability comparison between the two models. This would establish whether the EBM's $+5.93$ pp F1 advantage is a robust property of the dataset or sensitive to the benchmark's tuning state.

					\textit{R5---Transactional and behavioural feature integration.} Practitioner feedback identified the absence of behavioural or ongoing-suitability signals as the principal substantive gap of the present input schema. Future work, contingent on the data-access pre-condition established by Recommendation~R1, should extend the feature set with longitudinal transactional variables that approximate the behavioural signals a Philippine financial advisor would normally consult: recurring-deposit rate (deposits per month), withdrawal-to-deposit ratio, average idle-balance duration, balance-volatility coefficient over a rolling window, and median inter-transaction interval. These signals would complement the current snapshot preference variables (savings goal, risk tolerance, investment horizon) with a dynamic-behaviour layer aligned with the suitability-analysis convention codified in RA~11765 \cite{ref44}. Methodologically, the EBM's additive structure accommodates these features natively as additional shape functions without requiring a change of model class.

			\subsection{Closing Statement}
				This study began from the observation that Filipino fintech users are increasingly receiving automated financial-product recommendations under conditions of limited financial literacy \cite{ref4}, \cite{ref43}, nascent trust in AI-driven systems, and scarce access to transparent explanations of how and why a given product was recommended. To address this problem, the study designed and implemented FinXplain, an end-to-end explainable recommendation prototype grounded in Philippine statistical reality: a synthetic dataset of 5{,}000 profiles parameterized from PSA FIES \cite{ref40}, PSA LFS \cite{ref41}, the PSA 2020 Census \cite{ref39}, and the BSP Financial Inclusion Survey \cite{ref35}; a rule-based multi-factor labeling system traced to the official product documentation of GCash, Maya, and BPI \cite{ref6}, \cite{ref37}; an Explainable Boosting Machine model that achieved a weighted F1 of $0.8711$ and a weighted OvR AUC-ROC of $0.9761$, outperforming an untuned XGBoost benchmark by $+5.93$ and $+0.90$ percentage points respectively (Table~\ref{tab:threshold}); and a working Next.js and FastAPI web prototype that integrates EBM native explanations, plain-language reasons, and counterfactual suggestions into a single accessible interface. The key finding, in one sentence, is that a glass-box EBM can match---and on this dataset exceed---the predictive performance of a black-box XGBoost benchmark while preserving the native interpretability demanded by the Philippine Financial Products and Services Consumer Protection Act \cite{ref44} and by the broader theoretical case against post-hoc-only explanation in high-stakes domains \cite{ref51}. As the Philippine fintech ecosystem continues to grow and as the BSP advances its financial-inclusion and digital-payments agenda \cite{ref35}, \cite{ref75}, tools that help ordinary Filipinos understand not only which product is recommended for them but why---and what they can do to become eligible for other options---represent an important direction for responsible AI in financial services. This paper's specific contribution to that broader program is a working demonstration that, on a Philippine-contextualized synthetic fintech benchmark and under a regulation-aligned design, an inherently interpretable model combined with a multi-modal explanation layer is an implementable and competitive alternative to a post-hoc-explained black box; the next steps---real-data validation, usability testing with Filipino users---are explicitly identified in the preceding recommendations.

		%References command
		%Parameter is the references bib file without extension
		\makereferences{references}
			
	\end{mainmatter}
\end{document}
